\subsection{The algebra actually written on 155/002}

Throughout this subsection, $k$ is a field and $L/k$ is finite. All tensor
categories and tensor functors are symmetric and unital. We make this convention
explicit: without symmetry the commutativity of $U$ in the manuscript is not
the appropriate conclusion. A pro-vector space means the inverse system dual
to an ordinary vector space; completed tensor products have this meaning.
All ideals in such a space are closed. Finite-dimensional examples need no
completion.

\begin{saresult}{B.001}{Proved}{The idealizer is the Takeuchi product}
Let $A$ be a unital $k$-algebra containing $L$ by $s:L\longrightarrow A$.
Put $D=A\otimes_k A$ and let $J$ be the right ideal generated by
$s(\lambda)\otimes1-1\otimes s(\lambda)$, for $\lambda\in L$.
The quotient $D/J$ is the tensor product of the two left $L$-modules $A$.
The image of the left idealizer $\widetilde J=\{d:dJ\subset J\}$ is
\[
 A\mathbin{\times_L}A=
 \left\{\sum_i a_i\otimes_L b_i:
 \sum_i a_i s(\lambda)\otimes_L b_i
 =\sum_i a_i\otimes_L b_i s(\lambda)
 \text{ for every }\lambda\in L\right\}.
\]
Multiplication is induced by factorwise multiplication in $D$. Thus
$\widetilde J/J\simeq A\times_L A$ as algebras. The same statement holds
for the completed products of 155/002.
\end{saresult}
\begin{proof}
The defining relations for $D/J$ are
$s(\lambda)a\otimes b=a\otimes s(\lambda)b$, exactly the balancing
relations of the indicated tensor product. An element $d$ idealizes $J$
if and only if $d(s(\lambda)\otimes1-1\otimes s(\lambda))$ lies in $J$
for every $\lambda$, because these differences generate $J$ on the right.
This is the displayed condition. The idealizer is a subalgebra, and $J$
is a two-sided ideal in it: right multiplication preserves $J$ by its
definition, and left multiplication by an idealizer element preserves it
by the idealizer condition. This proves the algebra assertion, without
incorrectly assigning an algebra structure to all of $A\otimes_L A$.
The completed assertion follows by the same closed-relations argument,
or by taking inverse limits of the finite linear quotients.
\end{proof}

\begin{saresult}{B.002}{Refuted}{The displayed conditions on 155/002 omit a unit constraint}
Even for $L=k=\mathbb F_2$ and finite-dimensional $U$, the data 1), 2),
(i), and (ii) on 155/002 do not imply the asserted tensor-category
reconstruction. The missing constraint is the character identity for
the coalgebra counit. In the neutral case it says that this counit is
an algebra homomorphism. Consequently R7 and L6 require this correction.
\end{saresult}
\begin{proof}
Let $S=\{a,b,z\}$, with $a^2=a$, $b^2=b$, $ab=ba=z$, and $zs=sz=z$
for every $s\in S$. This is a commutative associative semigroup with no
identity. Let $U=kS$ be its semigroup algebra. It is nevertheless a
unital algebra: its identity is $e_a+e_b-e_z$. The map
$\varepsilon_U:U\longrightarrow k$, $e_s\longmapsto1$, is a unital
algebra homomorphism. Take $A=U^\vee=k^S$, with pointwise multiplication
and identity the constant function $1$, which is precisely
$\varepsilon_U$. Condition (i) is the ordinary inclusion of scalar
functions. The transpose of the multiplication of $U$ is
\[
 \Delta:A\longrightarrow A\otimes_k A=k^{S\times S},
 \qquad \Delta(f)(s,t)=f(st).
\]
It is a unital algebra homomorphism. Since $J=0$, this verifies (ii)
including its strongest usual unital interpretation. Coassociativity
and cocommutativity follow from the semigroup law. The coalgebra counit
dual to $1_U$ is
\[
 e:A\longrightarrow k,\qquad e(f)=f(a)+f(b)-f(z).
\]
For the indicator functions $p_a,p_b$, we have
$p_ap_b=0$ and $e(p_a)=e(p_b)=1$; hence $e$ is not multiplicative.
All these calculations remain valid in characteristic two.

There is also an obstruction independent of terminology. Finite
$A$-modules have three one-dimensional simples $V_s$, and the proposed
tensor product satisfies $V_s\otimes V_t=V_{st}$. A unit whose image
under the proposed fibre functor is $k$ must be some $V_s$. It would
require $st=t$ for all $t$, but $S$ has no such $s$. Thus the proposed
category has no permitted tensor unit. The attached program
\texttt{check\_B\_finite.py} exhausts all eight elements of $U$ over
$\mathbb F_2$, all $512$ associativity triples, all $64$ augmentation
pairs, and all $64$ coproduct-multiplicativity pairs. It also checks both
counit identities and all three possible one-dimensional units.
\end{proof}

\begin{saresult}{B.003}{Proved under stated hypotheses}{Corrected pro-bialgebroid reconstruction}
In the manuscript data, assume that $A$ is a pseudocompact $k$-algebra
and that the multiplication is continuous. Write $e:A\longrightarrow L$
for the transpose of $1_U$. Add
\begin{equation}\label{eq:B-character}
 e(1)=1,\qquad
 e(ab)=e\bigl(a\,s(e(b))\bigr)\quad(a,b\in A).
\end{equation}
The first equality already follows if the stated augmentation of $U$
is unital. With this addition, the displayed data are exactly the data
of a continuous cocommutative left $L$-bialgebroid with equal source
and target $s$. They give a symmetric abelian tensor category of
continuous $A$-modules finite-dimensional over $L$, with exact faithful
forgetful functor to $\operatorname{Modf}(L)$. Conversely every
essentially small abelian $k$-linear symmetric tensor category with
such a fibre functor reconstructs these data. Both reconstructions are
inverse up to an equivalence respecting the fibre functors and the
induced isomorphism of the pro-algebras.
\end{saresult}
\begin{proof}
For the forward direction, the tensor product of two modules has
underlying space $M\otimes_L N$ and action
\[
 a(m\otimes n)=\sum a_{(1)}m\otimes a_{(2)}n.
\]
The sum is interpreted in a finite quotient when acting on $M,N$.
B.001 proves that balancing over $L$ is respected. Multiplicativity
of $\Delta$ proves the module law; coassociativity gives associativity
of tensor products; cocommutativity gives the usual symmetry. The
tensor unit is $L$, with action
\[
 a\cdot\lambda=e(a s(\lambda)).
\]
Equation \eqref{eq:B-character} says exactly that this is an $A$-action:
$(ab)\cdot\lambda=a\cdot(b\cdot\lambda)$. The ordinary coalgebra
counit identities, applied to $\Delta(a s(\lambda))$, show that
$L\otimes_L M\longrightarrow M$ and $M\otimes_L L\longrightarrow M$
are $A$-linear. The action of $s(\lambda)$ is ordinary scalar
multiplication, since $e$ is left $L$-linear. All coherence maps are
the usual linear ones. Kernels and cokernels retain continuous
$A$-actions and their finite dimension; the forgetful functor is exact
and faithful. Tensoring is exact because $L$ is a field.

For reconstruction, put $W=\operatorname{Res}_{L/k}F$ and give
$A=\operatorname{Nat}_k(W,W)$ its finite topology. The category is
locally finite over $k$: its Hom spaces inject into finite-dimensional
linear Hom spaces, and every strictly increasing chain of subobjects
strictly increases $\dim_LF(X)$. The coalgebra reconstruction theorem
in Deligne \cite[6.1--6.2]{Deligne1990}, applied to $W$, therefore
identifies the category with the finite-dimensional continuous
$A$-modules. This is an application of the cited theorem, with all its
finiteness hypotheses checked, rather than an assertion that an
arbitrary faithful functor suffices. Scalar multiplication on $F(X)$
gives $s:L\longrightarrow A$.

For completeness, the algebraic passage from the tensor product to
$\Delta$ can be checked directly after this equivalence. A natural
endomorphism of the bifunctor $(M,N)\longmapsto M\otimes_L N$ is
determined on pairs of finite regular quotients $A/I,A/J$ by its value
at $1\otimes1$. Naturality under the maps $A/I\longrightarrow M$,
$a\longmapsto am$, and their analogues for $N$, determines all its
values. Naturality under right multiplication by $s(\lambda)$ is
exactly the Takeuchi condition in B.001. Compatible values on the
regular quotients are therefore precisely the elements of the
completed Takeuchi product. The action of each $a\in A$ on the
tensor-product object yields $\Delta(a)$. Associativity, symmetry,
and the action law give the coproduct axioms. The action on the tensor
unit gives $e(a)=a_1(1)$ and hence \eqref{eq:B-character}. Dualizing
the coproduct gives the commutative unital $L$-algebra $U$; dualizing
$1_A$ gives its augmentation. Finally the same regular-quotient
argument identifies the natural endomorphism algebra of the
forgetful functor on the reconstructed category with $A$ itself.
This proves both inverse comparisons. The categorical equivalence is
also the finite-field-base case of Hai
\cite[Theorem 2.1.2 and Corollary 2.2.5]{Hai2008}, with the manuscript's
pro-dual convention in place of his coring convention.
\end{proof}

\begin{saresult}{B.004}{Known}{The precise answer to 155/004, including inseparability}
Assume that the category in B.003 is rigid and that
$k\simeq\operatorname{End}(1)$. Its scalar extension to $L$ has an exact
faithful fibre functor induced by $F$. In the module description this
is the restricted diagonal functor on finite
$(A\widehat\otimes_k L)$-modules written on 155/004. It is exact and
faithful for every finite field extension $L/k$, including inseparable
ones. By contrast, the diagonal functor on all
$(L\otimes_k L)$-modules need not be exact, and is not generally
faithful even for a nontrivial separable extension.
\end{saresult}
\begin{proof}
Deligne \cite[Theorem 1.12, proved in 6.8 and 6.14--6.15]{Deligne1990}
identifies the rigid category with representations of the affine
groupoid $G\rightrightarrows\operatorname{Spec}L$, whose map to
$\operatorname{Spec}(L\otimes_k L)$ is faithfully flat. After extension
of the ground field to $L$, the object scheme has its diagonal
$L$-point. Restriction of an equivariant vector bundle to this point
is exact because the underlying modules in the representation
category are locally free. It is faithful because the groupoid is
transitive: after the faithfully flat base change parametrizing
arrows from that point to all objects, every equivariant map is
transported from its value at the point. A map vanishing there thus
vanishes after a faithfully flat base change and hence vanishes.
The groupoid and scalar-extension descriptions agree by Deligne's
scalar-extension construction in \S5 and its tensor-product
compatibility in 5.18--5.21. On modules the restriction is precisely
$M'\longmapsto M'\otimes_{L\otimes_k L}L$.

The unrestricted assertions are elementary. If
$L=k(u)$, $u^p=t\in k$ is purely inseparable of degree $p$, then
$L\otimes_k L\simeq L[\delta]/(\delta^p)$ and the diagonal quotient
is $\delta\mapsto0$. Tensoring the injection
$(\delta^{p-1})\hookrightarrow L[\delta]/(\delta^p)$ with $L$ gives
a nonzero source and a zero map. Thus it is not exact. For a nontrivial
finite separable extension, the diagonal component splits off as an
idempotent summand of $L\otimes_k L$; the other nonzero summand is
killed, so faithfulness fails. Neither test is a counterexample to
the restricted functor: arbitrary modules over the scalar subring
need not carry the required $A\otimes_k L$-action.
\end{proof}

\begin{saresult}{B.005}{Proved}{What completes the duality calculation on 155/006--010}
Choose a nonzero $k$-linear functional $\ell:L\longrightarrow k$ and
write $U=\operatorname{Hom}^{\mathrm{cont}}_L(A,L)$, where $A$ is a left
$L$-module by $s$. Give $U$ the left $A$-action
$(a u)(b)=u(ba)$. Then
\[
 \operatorname{Hom}_A(M,U)\simeq\operatorname{Hom}_L(M,L),
 \qquad \phi\longmapsto\bigl(m\longmapsto\phi(m)(1)\bigr),
 \qquad \operatorname{End}_A(U)\simeq A^{\mathrm{op}}.
\]
The last isomorphism uses $\ell$ when $L\ne k$. A rigid tensor
structure supplies a continuous anti-isomorphism
$a_\ell:A\longrightarrow A^{\mathrm{op}}$, fixing $s(L)$, and hence a
second, commuting $A$-action on $U$. Evaluation and coevaluation are
\[
 \operatorname{ev}_M(m\otimes\phi)=\phi(m)(1),\qquad
 \operatorname{coev}_M(1)=\sum_i m_i\otimes\phi_i,
\]
where $(m_i), (\phi_i)$ are dual $L$-bases under the first isomorphism.
Conversely a continuous anti-isomorphism fixing $s(L)$ gives rigidity
exactly when these two maps are $A$-linear for every finite continuous
$A$-module $M$. Their two triangle identities always hold as linear
maps, and therefore hold in the module category when the maps are
$A$-linear. For pro-$A$, the manuscript's instruction to test only
$M=A$ must be replaced by tests on finite modules; $A$ need not itself
be an object of the category.
\end{saresult}
\begin{proof}
Given $f\in\operatorname{Hom}_L(M,L)$, set
$\phi_f(m)(b)=f(bm)$. This is left $L$-linear in $b$, continuous, and
$A$-linear in $m$. It inverts the displayed evaluation-at-$1$ map.
The bilinear form $(x,y)\longmapsto\ell(xy)$ on the field $L$ is
nondegenerate: for $x\ne0$, choose $z$ with $\ell(z)\ne0$ and put
$y=x^{-1}z$. Consequently
\[
 U\longrightarrow\operatorname{Hom}^{\mathrm{cont}}_k(A,k),
 \qquad u\longmapsto\ell\circ u
\]
is an isomorphism. Its inverse sends $g$ to the unique $u$ satisfying
$\ell(\lambda u(b))=g(s(\lambda)b)$ for all $\lambda\in L$.
The map intertwines the coinduced $A$-actions. Continuous duality
therefore gives
\[
 \operatorname{End}_A(U)
 \simeq\operatorname{End}_{A^{\mathrm{op}}}^{\mathrm{cont}}(A)^{\mathrm{op}}
 \simeq A^{\mathrm{op}},
\]
because a continuous endomorphism of the right regular module is
left multiplication by its value at $1$. Under this identification,
scalar multiplication on the values of $u$ corresponds to $s(L)$.

Suppose now that $D$ is the rigid duality. The fibre functor identifies
$F(DM)$ with the $L$-dual of $F(M)$. Pair it with $F(M)$ by
$B_M(\phi,m)=\ell(\phi(m))$. For $a\in A$, define
$a_\ell(a)_M$ as the adjoint of $a_{DM}$ under this nondegenerate
$k$-pairing. Naturality of the pairing makes this a natural
endomorphism of $F$ as a $k$-valued functor. Transposition reverses
products, and the bidual isomorphism gives the inverse map.
Multiplication by $\lambda\in L$ is self-adjoint, so $a_\ell$ fixes
$s(L)$. This is the required anti-isomorphism. The above adjunction
then identifies the dual module with $\operatorname{Hom}_A(M,U)$.
Its evaluation and coevaluation are the images of the corresponding
morphisms of the given rigid category, and hence are $A$-linear.

Conversely the stated anti-isomorphism defines a contravariant
equivalence on finite modules through $\operatorname{Hom}_A(-,U)$;
on underlying spaces it is exact ordinary $L$-duality. If evaluation
and coevaluation are $A$-linear, they give duality in the tensor
category. Explicitly the first triangle sends
$m$ to $\sum_i m_i\phi_i(m)=m$, and the second sends
$\phi$ to $\sum_i\phi(m_i)\phi_i=\phi$. These computations also prove
basis independence. A fully explicit criterion for $A$-linearity is
\[
 \sum\operatorname{ev}_M(a_{(1)}m\otimes a_{(2)}\phi)
 =a\cdot\operatorname{ev}_M(m\otimes\phi),
\]
\[
 \sum_i\sum a_{(1)}m_i\otimes a_{(2)}\phi_i
 =e(a)\sum_i m_i\otimes\phi_i.
\]
Here $a\cdot\lambda=e(a s(\lambda))$, so both equations are typed
over $L$. These are the evaluation and coevaluation constraints
missing from a bare anti-homomorphism. A finite separable extension
has the canonical permissible choice $\ell=\operatorname{Tr}_{L/k}$;
for a purely inseparable extension the trace is zero and cannot be
used for this purpose. We assert existence of a choice, not a
canonical inseparable trace.
\end{proof}

\subsection{The divided-power analyser and the missing biring data}

\begin{saresult}{B.006}{Proved}{The nonunital divided-power analyser}
Let
\[
 \Omega_{\mathrm{DP}}=\bigoplus_{n\ge1}\mathbb Z\gamma_n,
 \qquad \gamma_n=T^n/n!\in\mathbb Q[T].
\]
With ordinary polynomial addition, multiplication, and composition,
this is a nonunital analyser satisfying 153/006, axioms 1) and 2),
with composition unit $T=\gamma_1$ and constant subring zero. Its
natural domain of action is any divided-power ideal $I$ in a
commutative unital ring, by
$\left(\sum a_n\gamma_n\right)(x)=\sum a_n\gamma_n^I(x)$.
It does not satisfy the optional axiom 3a) literally printed on
153/006. The correct nonunital coaddition has its target in the
augmentation ideal of
$\mathbb Z[\gamma_1,\gamma_2,\ldots]\otimes
 \mathbb Z[\gamma_1,\gamma_2,\ldots]$, not in
$\Omega_{\mathrm{DP}}\otimes\Omega_{\mathrm{DP}}$ alone.
\end{saresult}
\begin{proof}
The product is
$\gamma_m\gamma_n=\binom{m+n}{m}\gamma_{m+n}$.
For $G=\sum_{i=1}^d a_i\gamma_i$ and $m\ge1$, the coefficient of
$\gamma_N$ in $\gamma_m(G)$ is
\[
 \sum_{\substack{r_1+\cdots+r_d=m\\\sum i r_i=N}}
 \frac{N!}{\prod_i r_i!(i!)^{r_i}}\prod_i a_i^{r_i}.
\]
The factorial quotient counts partitions of a labelled $N$-element
set into $r_i$ unordered blocks of size $i$ for each $i$, and is an
integer. Thus composition closes in $\Omega_{\mathrm{DP}}$.
Its associativity and unit follow from the injective embedding into
$\mathbb Q[T]$, and substitution in the outer variable is a
homomorphism of nonunital rings. A constant under composition must
equal its value on $0$, which is $0$ here. The divided-power axioms
give exactly the displayed product and substitution identities on
any divided-power ideal, so the stated action is defined even when
the ambient ring has torsion; division in that ring is neither
used nor asserted.

In any nonzero analyser with zero constant subring, optional 3a)
is impossible: put $G''=0$ there. Each $F_i''\circ0$ is a constant,
hence zero, so the right side vanishes, whereas for $F=T$ the left
side is $G'$. In the divided-power case the valid formula is
\[
 \gamma_n(x+y)=\gamma_n(x)+\gamma_n(y)
   +\sum_{i=1}^{n-1}\gamma_i(x)\gamma_{n-i}(y).
\]
The first two terms explain the required unitalization of the target.
For multiplication the formula
$\gamma_n(xy)=n!\gamma_n(x)\gamma_n(y)$ already lies in the
nonunital tensor square. The attached program verifies the pure
composition coefficients for $1\le m,n\le12$ and $461$ Bell
coefficients with at most six blocks of five possible sizes; the
set-partition argument, not this finite range, proves integrality
for all indices.
\end{proof}

\begin{saresult}{B.007}{Proved}{Exact extra data turning a unital analyser into a plethory}
Let $P$ be a unital analyser and let $k$ be a specified unital subring
of its constant ring. A compatible $k$-plethory structure is exactly
the following additional data and conditions: a $k$--$k$ biring
structure on the $k$-algebra $P$, with coaddition and comultiplication
realizing the finite decompositions in 3a) and 3b), with
$\beta(c)(F)=F\circ c$; and the requirement that the induced map
$P\odot_k P\longrightarrow P$, $F\odot G\longmapsto F\circ G$,
and $k[T]\longrightarrow P$, $T\longmapsto T$, are biring morphisms.
Here a biring includes additive and multiplicative counits, additive
coinverse, coassociativity, cocommutativity, codistributivity, and the
scalar map $\beta:k\longrightarrow\operatorname{Hom}_{k\text{-alg}}(P,k)$.
None of these coalgebra identities is encoded by unspecified
existence of decompositions alone.

There is a useful sufficient criterion with no choice of additional
operations: if the constant ring is a field $k$, $P$ acts faithfully
on $k$, and 1)--3) hold with the unitality convention of 153/006,
then there is a unique compatible $k$-plethory structure.
\end{saresult}
\begin{proof}
The first assertion is the explicit form of a monoid object in the
category of $k$--$k$ birings; it is the definition of Borger--Wieland
\cite[1.1, 1.3, 1.8]{BorgerWieland2005}. To check that it applies
to the given analyser, the relations defining $P\odot_k P$ are
linearity and multiplicativity in the first variable, the two
decomposition formulas in the second variable, and
$F\odot c=\beta(c)(F)$. These hold exactly by the specified
compatibility, so composition descends. The analyser monoid
identities give associativity and the two unit identities; the
remaining requirement is precisely that they are maps in the
biring category. Conversely every plethory supplies all these
data and its elementwise operations satisfy the analyser axioms.

For the sufficient criterion, identify $P$ with a $k$-subalgebra of
$k^k$ by faithful evaluation on constants. For every positive integer
$r$, the evaluation map
\[
 P^{\otimes_k r}\longrightarrow k^{k^r},\qquad
 F_1\otimes\cdots\otimes F_r\longmapsto
 \bigl((x_1,\ldots,x_r)\mapsto\prod_i F_i(x_i)\bigr)
\]
is injective. To see this for two factors, take a finite linearly
independent collection of functions in the first factor. The
evaluation vectors span its full linear dual, since a nonzero
linear combination cannot vanish at every point. Evaluating a
tensor at sufficiently many first coordinates therefore recovers
each coefficient function in the second factor. Induction proves
the assertion for every $r$.

The decompositions in 3a), 3b), evaluated on constants, now define
unique $\Delta_+F,\Delta_\times F\in P\otimes_k P$. Every biring
identity can be tested by the injective evaluation map with two,
three, or four variables, where it reduces respectively to an
ordinary ring identity in $k$. The counits are $F\mapsto F(0)$
and $F\mapsto F(1)$, the coinverse is $F\mapsto F\circ(-T)$,
and $\beta(c)$ is evaluation at $c$. The maps $\Delta_+$ and
$\Delta_\times$ are $k$-algebra maps by the same test. The
decomposition relations for composition already hold on all
elements of $P$ by hypothesis; hence composition descends to
$P\odot_k P$. Its compatibility with coaddition follows, on a
generator $F\odot G$, by evaluating both sides at $x,y\in k$:
both give $F(G(x+y))$. The multiplicative test gives $F(G(xy))$;
the counit, scalar, and coinverse tests are the corresponding
one-variable identities. Injectivity proves the biring identities.
The unit $k[T]\longrightarrow P$ has the required co-operations
because $T$ evaluates as the identity function. This completes the
construction and proves uniqueness.
\end{proof}

\begin{saresult}{B.008}{Refuted}{Freeness is insufficient for the ring extension on 153/004}
After the manuscript replaces a fraction-field inclusion by an
arbitrary inclusion $R\subset S$ of commutative rings, the marginal
claim ``OK if $\Omega$ is free'' does not ensure well-determined
coaddition or comultiplication tensors. It fails even when $R,S$
are fields. The correction ``$R$ is a domain'' in the earlier study
does not suffice either.
\end{saresult}
\begin{proof}
Take $R=\mathbb F_2$, $S=\mathbb F_4=\mathbb F_2[\alpha]$ with
$\alpha^2+\alpha+1=0$, and
\[
 \Omega=\{F\in S[T]:F(R)\subset R\},\qquad q=T^2-T.
\]
It is a free $R$-module, since $R$ is a field. Both $q$ and
$\alpha q$ belong to $\Omega$ and are linearly independent over
$R$. Thus
\[
  z=(\alpha q)\otimes q-q\otimes(\alpha q)
\]
is nonzero in $\Omega\otimes_R\Omega$. For example, extend
$q,\alpha q$ to an $R$-basis and apply the tensor product of their
two coordinate functionals. Nevertheless its image in $S[X,Y]$
is zero. A decomposition of the zero polynomial already has the
two distinct preimages $0,z$. Hence tensors in the claimed
expansions are not well-determined. The attached computation
exhausts the $16$ elements of the tensor square of
$\operatorname{span}_R(q,\alpha q)$: its polynomial image has four
elements and its kernel has four elements, including the displayed
witness. There is no finite-degree inference here: the witness in
this subspace remains nonzero in the full tensor product because
vector-space inclusions split.
\end{proof}

\begin{saresult}{B.009}{Proved}{The precise tensor criterion and the valid free-module theorem}
For an inclusion $R\subset S$, let
$\Omega=\{F\in S[T]:F(R)\subset R\}$ and let
\[
 \theta_2:\Omega\otimes_R\Omega\longrightarrow S[X,Y],
 \qquad G\otimes H\longmapsto G(X)H(Y).
\]
The polynomial identities on 153/004 determine one and only one
tensor for every $F$ if and only if $\theta_2$ is injective and both
$F(X+Y)$ and $F(XY)$ belong to its image for every $F\in\Omega$.
For a domain $R$ with fraction field $S$, freeness of $\Omega$ over
$R$ implies these conditions and, more strongly, gives isomorphisms
\[
 \theta_n:\Omega^{\otimes_R n}\xrightarrow{\sim}
 \{P\in S[T_1,\ldots,T_n]:P(R^n)\subset R\}
 \quad(n\ge1).
\]
In this case the coaddition, comultiplication, and ordinary
composition give a unique compatible $R$-plethory structure.
\end{saresult}
\begin{proof}
Existence of the requested tensors is exactly the image condition.
Uniqueness for $F=0$ already forces $\ker\theta_2=0$, and injectivity
then gives uniqueness for every $F$. This proves the exact criterion
for an arbitrary inclusion, including rings with zero divisors.

Now let $S=\operatorname{Frac}(R)$. If $R$ is finite it is a field,
$S=R$, and $\Omega=R[T]$, so the conclusion is immediate. Otherwise
$R$ is infinite. Let $\mathcal B$ be an $R$-basis of $\Omega$.
Since $R[T]\subset\Omega\subset S[T]$, localization gives
$S\otimes_R\Omega=S[T]$, and $\mathcal B$ is an $S$-basis of
$S[T]$. A polynomial $P(X,Y)$ integer-valued on $R^2$ has a unique
finite expansion $\sum_i b_i(X)h_i(Y)$ with $b_i\in\mathcal B$ and
$h_i\in S[Y]$. For every $y\in R$, the polynomial $P(X,y)$ belongs
to $\Omega$. Uniqueness of the $S$-basis expansion then implies
$h_i(y)\in R$. Therefore $h_i\in\Omega$, which proves surjectivity
of $\theta_2$. Injectivity follows because the tensor products of
the basis vectors remain linearly independent after localization.
The same argument in successive variables gives every $\theta_n$.

Transport the polynomial substitutions $T\mapsto X+Y$ and
$T\mapsto XY$ across $\theta_2$. The biring identities hold in
$S[X_1,\ldots,X_n]$, so injectivity of $\theta_n$ proves them in
the tensor powers. Evaluation at constants, negation, and the
inclusion $R[T]\subset\Omega$ provide the remaining biring maps.
Composition closes because an integer-valued polynomial sends
$R$ into $R$, and its compatibility identities are ordinary
substitution identities, again transported across the injective
$\theta_n$. Hence it defines the plethory structure. This proves
the manuscript's free-module argument under the fraction-field
hypothesis it originally had, but not under the later, weaker
hypothesis. Elliott's Theorem 4 and Proposition 13
\cite{Elliott2010} place this criterion in the more general
theory of weak polynomial compositeness. Its questions record the
state of that 2010 paper, not a claim that they remain open today;
B.016 supplies a domain where even multiplicative existence fails.
\end{proof}

\begin{saresult}{B.010}{Proved}{Faithfulness of evaluation on constants}
For any analyser $\Omega$, evaluation on its constant subring
$\Omega_0$ is injective exactly when the constants separate
elements of $\Omega$ under composition. For the polynomial analyser
of an inclusion $R\subset S$, this has the following explicit
coefficient criterion:
\[
 \bigcap_{r\in R}\ker\bigl(S^{d+1}\longrightarrow S,
 (a_0,\ldots,a_d)\longmapsto\sum_{i=0}^d a_i r^i\bigr)=0
 \quad\hbox{for every }d\ge0.
\]
In particular, if $S$ is a domain it is injective exactly when
$R$ is infinite. More generally it is injective if for every $d$
there are $r_0,\ldots,r_d\in R$ such that
$\prod_{i<j}(r_j-r_i)$ is a non-zero-divisor in $S$.
The unqualified desirability statement on 153/008 is not a theorem:
$\mathbb F_q[T]$ and the divided-power analyser B.006 both give
nonfaithful examples.
\end{saresult}
\begin{proof}
The first assertion is the kernel definition of the indicated
representation, stated to distinguish an analyser from its
possibly nonfaithful action on constants. In the polynomial case,
the constants are exactly $R$: a polynomial constant under
composition equals its composition with $0$, and therefore is
constant as a polynomial. Every polynomial vanishing on $R$ lies
in $\Omega$. Thus the kernel is the null-polynomial ideal
$\{F\in S[T]:F(r)=0\text{ for all }r\in R\}$, and collecting
coefficients by degree gives the displayed exact criterion.
A nonzero polynomial over a domain has at most its degree many
roots, proving the infinite-domain case. If $R$ is finite,
$\prod_{r\in R}(T-r)$ is a nonzero monic null polynomial, so
evaluation is not injective even if $S$ has zero divisors.
Under the stated non-zero-divisor condition, the Vandermonde
matrix on $r_0,\ldots,r_d$ has determinant
$\prod_{i<j}(r_j-r_i)$. Multiplication by its adjugate shows that
a vector in its kernel is annihilated by that determinant and
is therefore zero. This proves the sufficient condition in
general. The specific null polynomial over $\mathbb F_q$ is
$T^q-T$. B.006 has zero constant ring but contains the nonzero
element $T$, so its action on constants is also nonfaithful.
\end{proof}

\begin{saresult}{B.011}{Refuted}{A finite unital analyser satisfying all three axioms need not be a plethory}
The $64$-element algebra
\[
 P=\mathbb F_2[T]/\bigl((T^2-T)^3\bigr)
\]
with ordinary polynomial substitution is a unital analyser
satisfying 153/006, 1)--3), including both optional decomposition
axioms. It has no compatible plethory structure. Thus the missing
biring data in B.007 cannot be inferred from those axioms.
\end{saresult}
\begin{proof}
Put $q=T^2-T$. For every $G\in\mathbb F_2[T]$, the polynomial
$G^2-G$ vanishes at $0,1$, so it is divisible by $q$.
Consequently $q(G)^3$ is divisible by $q^3$; changing an outer
polynomial modulo $q^3$ does not change its substituted value.
Changing the inner polynomial by $q^3H$ also does not change the
result modulo $q^3$, by the factor theorem applied to
$F(X)-F(Y)$. Thus composition descends. Its associativity, unit
$T$, and the homomorphism property in the outer variable descend
from polynomials. The ring unit is the constant $1$.

For a representative $F(T)=\sum_{n=0}^5 a_nT^n$, expand
$F(X+Y)$ by the binomial theorem and $F(XY)$ term by term.
These expansions are finite sums of products of one-variable
polynomials. Substituting arbitrary elements $G,H\in P$ gives
exactly axioms 3a), 3b). Their dependence on the choice of a
representative is immaterial to the existential axioms.

If a compatible plethory existed, its unit map
$\mathbb F_2[T]\longrightarrow P$ would be a biring map.
Therefore $\Delta_+(T)=T\otimes1+1\otimes T$ and
$\Delta_+(q)=q\otimes1+1\otimes q$. Since $\Delta_+$ is an
algebra map and $q^3=0$ in $P$, we would obtain
\[
 0=\Delta_+(q^3)=q^2\otimes q+q\otimes q^2.
\]
This tensor is nonzero: $q,q^2$ are linearly independent in the
six-dimensional vector space $P$, and their tensor products are
independent. This is a contradiction. The attached program
exhaustively checks all $64^3=262144$ triples for composition
associativity, both homomorphism identities, and both optional
decomposition identities; it also computes the nonzero tensor.
\end{proof}

\subsection{Comparison with the primary reconstruction sources}
\begin{saresult}{B.012}{Known}{What cote 155 does and does not anticipate}
After adding B.003's unit constraint, the algebra on 155/002 is a
cocommutative equal-source-and-target instance of Takeuchi's
$\times_L$-bialgebra, presented as a pro-dual. Without that constraint,
it is not a bialgebroid, by B.002. The manuscript's reconstruction is a
special case of the general reconstruction theorems below, with a
separate completion step; its rigid case has the non-neutral groupoid
interpretation of Deligne. Thus L6 is valid as a qualified structural
comparison, and is false if it asserts that the literal displayed data
already satisfy the modern definition. No priority claim or claim of
independence follows from these comparisons.
\end{saresult}
\begin{proof}
The primary sources have been opened and the following exact portions
checked.

Takeuchi \cite[\S1, especially 1.1--1.2; Definition 4.5, pp.474--475]
{Takeuchi1977} defines the balancing-and-invariance product and requires
both the coproduct and the map to $\operatorname{End}_k(L)$ to be
algebra maps. Substituting the two equal copies of the commutative base
$L$ gives B.001. His map to $\operatorname{End}_k(L)$ is
$a\mapsto(\lambda\mapsto e(a s(\lambda)))$ in B.003, so its
multiplicativity is exactly the missing character equation. This identifies
the structure, including what the manuscript omitted; the definition
itself is not an abelian-category reconstruction theorem.

Schauenburg \cite[Theorem 5.1 and Remark 5.3; author's preprint pp.15--17]
{Schauenburg1998} proves that an algebra over $R\otimes R^{\mathrm{op}}$
has a $\times_R$-bialgebra structure precisely when its module category
has a monoidal structure with strictly monoidal underlying functor to
$R$-bimodules. Remark 5.3 removes strictness by transport of structure.
The proof reconstructs the coproduct from the regular module and the
counit from the action on the tensor unit. This is exactly B.003's
algebra-to-tensor-structure comparison before completion. For the
manuscript's pro-algebra, the regular module need not be finite, so one
must replace it by all finite regular quotients, as our proof does.
Schauenburg's theorem by itself does not reconstruct an arbitrary abelian
category from a faithful exact functor.

Hai \cite[Theorem 2.1.2, Theorem 2.2.4, Corollary 2.2.5, and
Theorem 2.2.6]{Hai2008} supplies that reconstruction by a coend: local
finiteness, exactness, faithfulness and finite projectivity give the
abelian equivalence; monoidality gives the bialgebroid; left and right
rigidity give the corresponding Hopf structures. B.003 checked local
finiteness, and finite-dimensional $L$-spaces are projective, so these
hypotheses hold here. His algebra is the coend acting by coactions;
the manuscript's $A$ is the continuous pro-dual acting by actions.
They must not be identified without this dualization or with a
finite-dimensional dual unless finiteness has actually been proved.

Finally Schauenburg \cite[\S3 and Theorems 5.1, 5.5]{Schauenburg2003}
identifies $\times_R$-bialgebras over a Frobenius-separable base with
weak bialgebras. A finite separable field extension admits that
Frobenius-separable structure; a purely inseparable extension does not.
Thus this particular weak-bialgebra reformulation has a separability
restriction, while B.003 and Deligne's transitive-groupoid theorem in
B.004 do not. None of these theorems makes the unrestricted diagonal
functor on all $L\otimes_kL$-modules faithful or exact. This separates
the three logically different uses of separability and completes the
requested comparison.
\end{proof}

\subsection{Further classifications and an explicit domain obstruction}
\begin{saresult}{B.013}{Proved}{Complete answer for a finite base ring}
Let $R\ne0$ be a finite commutative ring and let $R\subset S$ be a
unital inclusion. For the module $\Omega$ and map $\theta_2$ of B.009,
the following are equivalent: $\theta_2$ is injective; the identities
on 153/004 determine unique coaddition and comultiplication tensors;
and $S=R$. In the last case $\Omega=R[T]$ with its usual plethory
structure. This classifies arbitrary extensions of a finite base,
including extensions that are infinite or have zero divisors.
\end{saresult}
\begin{proof}
Set $P(T)=\prod_{r\in R}(T-r)$ and $d=|R|$. This is a monic
polynomial of degree $d$ vanishing on $R$. Division by $P$, valid over
any commutative coefficient ring because $P$ is monic, gives
\[
 \Omega=E\oplus P S[T],\qquad
 E=\{H\in S[T]:\deg H<d,\ H(R)\subset R\}.
\]
Indeed $F=H+PQ$ is integer-valued on $R$ exactly when $H$ is.
The summand $SP$ of $PS[T]$ is itself split: project $Q$ to its
constant coefficient. Thus $S\otimes_R S$ embeds as a direct summand
of $\Omega\otimes_R\Omega$ by $s\otimes t\mapsto(sP)\otimes(tP)$.
On this summand $\theta_2$ is
\[
 s\otimes t\longmapsto stP(X)P(Y).
\]
Injectivity of $\theta_2$ forces multiplication
$\mu:S\otimes_R S\to S$ to be injective. It is always surjective,
so it is an isomorphism.

Write $J$ for the Jacobson radical of $R$. It is nilpotent: every
$x\in J$ has repeating powers, say $x^a=x^{a+b}$, and $1-x^b$ is a
unit, whence $x^a=0$; a finite generating set of this nil ideal gives
$J^N=0$ by the pigeonhole principle. The Chinese remainder theorem
identifies $R/J$ with a finite product of finite fields. After base
change, multiplication for $T=S/JS$ over $R/J$ is still an
isomorphism. On a field factor $k$, either $T_k=0$, or choose a
$k$-linear functional $\ell:T_k\to k$ with $\ell(1)=1$.
For every $a\in T_k$, injectivity of multiplication gives
$a\otimes1=1\otimes a$. Applying $\ell\otimes1$ yields
$a=\ell(a)1$. Thus every nonzero factor equals $k$ and
$R/J\to S/JS$ is surjective. Consequently $S=R+JS$.
Successive substitution gives $S=R+J^nS$ for all $n$, and $J^N=0$
gives $S=R$. Conversely, if $S=R$ then $\Omega=R[T]$ and every
$\theta_n$ is the standard polynomial-ring isomorphism. B.009 now
proves the equivalence with the requested uniqueness statement.
This assertion concerns uniqueness of the polynomial tensor
expansions; it does not infer nonexistence of every conceivable
abstract biring structure merely from a nonzero tensor kernel.
\end{proof}

\begin{saresult}{B.014}{Proved}{All field bases and the rational-coefficient reduction}
Let $R$ be an infinite domain, $K=\operatorname{Frac}(R)$, and let
$R\subset S$ with $S$ torsion-free as an $R$-module. Inside
$K\otimes_R S$ define the overring
\[
 S_{\mathrm{rat}}=S\cap(K\cdot1).
\]
Then
\[
 \Omega(R,S)=\{F\in S_{\mathrm{rat}}[T]:F(R)\subset R\}.
\]
Moreover $\ker\theta_2$ is exactly the $R$-torsion submodule of
$\Omega\otimes_R\Omega$. In particular, for an infinite field base
$R$, every unital extension $R\subset S$ gives $\Omega=R[T]$ and
unique standard co-operations. Together with B.013, this is a
complete classification for every field base: an infinite field
allows every extension; a finite field allows unique polynomial
tensor expansions exactly for the identity extension.
\end{saresult}
\begin{proof}
Torsion-freeness embeds $S$ in $K\otimes_R S$. If $F$ has degree
at most $d$, choose distinct $r_0,\ldots,r_d\in R$. Its coefficient
vector is obtained from $(F(r_0),\ldots,F(r_d))\in R^{d+1}$ by the
inverse Vandermonde matrix over $K$. Therefore each coefficient
belongs to $K\cdot1$ as well as to $S$, proving the first equality.
The converse inclusion is immediate. Since
$R[T]\subset\Omega\subset K[T]$, localization identifies
$K\otimes_R\Omega$ with $K[T]$; after localization $\theta_2$ is the
injective standard map $K[X]\otimes_KK[Y]\to(K\otimes_RS)[X,Y]$.
Its kernel is therefore torsion. Conversely a torsion tensor maps
to zero because the target $S[X,Y]$ is $R$-torsion-free. This proves
the asserted equality. If $R=K$ is an infinite field, then
$S_{\mathrm{rat}}=R$ and every polynomial in $R[T]$ preserves $R$,
which proves the field assertion.
\end{proof}

\begin{saresult}{B.015}{Proved}{A two-variable difference obstruction}
Let $(R,\mathfrak m)$ be a local domain with fraction field $K$.
Suppose $u_n\in R$ have the following property: for every finite
set $\mathcal H\subset\operatorname{Int}(R)$ there is an $n$ with
$h(u_n)-h(0)\in\mathfrak m$ for every $h\in\mathcal H$.
If $g\in\operatorname{Int}(R)$ satisfies $g(0)=0$ and
$g(u_n^2)\notin\mathfrak m^2$ for every $n$, then
\[
 g(XY)\notin\operatorname{im}\bigl(
 \operatorname{Int}(R)\otimes_R\operatorname{Int}(R)
 \longrightarrow K[X,Y]\bigr).
\]
Consequently no compatible biring, and hence no compatible plethory,
can have its polynomial inclusion $R[T]\to\operatorname{Int}(R)$
as a biring morphism.
\end{saresult}
\begin{proof}
If $g(XY)=\sum_{i=1}^N f_i(X)h_i(Y)$, choose $u=u_n$ for the
finite collection of all the $f_i,h_i$. Take the difference between
the values at $(u,u),(u,0),(0,u),(0,0)$ with signs $+,-,-,+$.
On the right it is
\[
 \sum_i\bigl(f_i(u)-f_i(0)\bigr)
          \bigl(h_i(u)-h_i(0)\bigr)\in\mathfrak m^2.
\]
On the left it is $g(u^2)$, a contradiction. For the last assertion,
localize a putative multiplicative co-operation to $K$.
Because $K\otimes_R\operatorname{Int}(R)=K[T]$ and the unit is a
biring map, its localized value on $T$ is $XY$, so its value on
$g$ must be $g(XY)$. Its unlocalized value would be precisely a
tensor preimage whose nonexistence was just proved. This argument
requires neither injectivity of $\theta_2$ nor a general
classification theorem for integer-valued polynomial rings.
\end{proof}

\begin{saresult}{B.016}{Refuted}{Explicit one-dimensional local domains without comultiplication}
For every finite field $k=\mathbb F_q$, put
\[
 T=k[t]_{\{h:h(0)h(1)\ne0\}},\qquad m=t(t-1),\qquad
 R=k+mT\subset K=k(t).
\]
Then $R$ is a one-dimensional Noetherian local domain, but
$\operatorname{Int}(R)$ has no compatible comultiplication tensor
for a specific polynomial of degree $q^2$. In particular the
fraction-field hypothesis alone does not repair the unrestricted
claim on 153/004; freeness or another genuine tensor condition is
needed. The case $q=2$ was found in the opened primary preprint
\cite[Theorem 3.4, pp.7--8]{JiangEtAl2026}; the argument here treats
all finite fields and proves the final biring obstruction directly.
\end{saresult}
\begin{proof}
The two maximal ideals of $T$ are $tT$ and $(t-1)T$, with residue
field $k$. Their intersection is $M=mT$. Thus $R$ consists exactly
of those elements of $T$ with equal residues at the two points.
An element with nonzero common residue has its inverse in $T$ with
equal residues, hence its inverse lies in $R$. Therefore $R$ is
local with maximal ideal $M$ and $R/M=k$. It is a domain because
it is a subring of $k(t)$, and its fraction field is $k(t)$ since
$m,mt\in R$ and $t=(mt)/m$.

For the stated finiteness properties, let $A=k[m]_{(m)}$. The ring
$A[t]$ is free over $A$ on $1,t$, using $t^2=t+m$. If
$h(t)\in k[t]$ is nonzero at $0,1$, its product with $h(1-t)$
belongs to $k[m]$ and has nonzero constant term. Thus $h(t)$ is
invertible in $A[t]$, and $T=A\oplus At$. It follows that
\[
 R=A\oplus A(mt)\simeq
 A[v]/(v^2-mv-m^3).
\]
This is a finite $A$-algebra and is Noetherian. Any nonzero prime
of $R$ meets $A$ nontrivially: multiplication by a nonzero element
of the prime on the free rank-two $A$-module $R$ has nonzero
determinant, and its characteristic equation puts that nonzero
determinant in the prime. Hence every nonzero prime contracts to
$(m)$ and is the unique maximal ideal $M$. The chain
$(0)\subsetneq M$ proves that the dimension is one.

Let $v_0,v_1$ be the valuations at $t=0,1$, and set
\[
 u_n=t(t-1)^{n+1},\qquad
 a(X)=\frac{X^q-X}{m},\qquad g(X)=a(X)^q-a(X).
\]
For $x\in R$, $x^q-x\in M$, so $a(x)\in T$. Both residues of
$a(x)$ belong to $\mathbb F_q$, whence $g(x)\in M\subset R$.
Thus $g\in\operatorname{Int}(R)$ and $g(0)=0$. Since
$v_0(u_n)=1$,
\[
 v_0\bigl(a(u_n^2)\bigr)
 =v_0\bigl((u_n^{2q}-u_n^2)/m\bigr)=1,
 \qquad v_0\bigl(g(u_n^2)\bigr)=1.
\]
Every element of $M^2=m^2T$ has $v_0\ge2$, so the second hypothesis
of B.015 holds. For a finite collection of polynomials
$h(X)=\sum_r c_{h,r}X^r$, choose $n$ so large that
$v_1(c_{h,r})+r(n+1)>0$ for every nonconstant coefficient in the
collection. Then $h(u_n)-h(0)\in(t-1)T$; both values lie in $R$,
so their difference belongs to $(t-1)T\cap R=M$. This is the first
hypothesis of B.015, which proves the obstruction.

For completeness, $g$ has leading term $m^{-q}X^{q^2}$ and hence
has degree exactly $q^2$. The attached symbolic check verifies the
valuation formulas for $q=2,3,4,5,7,8,9$ and $1\le n\le12$;
the valuation argument above proves them for every prime power
and every positive integer $n$.
\end{proof}

\begin{saresult}{B.017}{Proved}{Which base rings work for every extension}
For a nonzero commutative ring $R$, the following are equivalent:
\begin{enumerate}
\item every residue field $R/\mathfrak m$ at a maximal ideal is infinite;
\item for every unital inclusion $R\subset S$, every polynomial in
$S[T]$ preserving $R$ already belongs to $R[T]$;
\item for every such inclusion, the identities of 153/004 determine
unique coaddition and comultiplication tensors.
\end{enumerate}
Thus no structural classification gate remains for bases all of whose
residue fields are infinite. This includes every ring containing an
infinite field. If $R$ has even one finite residue field, the square-zero
extension constructed below already violates uniqueness.
\end{saresult}
\begin{proof}
Assume condition 1 and let $F=\sum_{i=0}^d a_iT^i\in S[T]$
satisfy $F(R)\subset R$. Fix a maximal ideal $\mathfrak m$ of $R$.
Choose $d+1$ elements of $R$ with distinct images in the infinite
field $R/\mathfrak m$. Their Vandermonde determinant is outside
$\mathfrak m$, so the inverse matrix exists over $R_{\mathfrak m}$.
Apply it to the evaluations of $F$. Each $a_i$ has image in
$R_{\mathfrak m}\subset S_{\mathfrak m}$, so the class of $a_i$ in
the $R$-module $S/R$ vanishes after localization at every maximal
ideal. Such a class is zero: if a nonzero element $x$ of a module
had this property, its proper annihilator would be contained in
some maximal ideal $\mathfrak m$, whereas $x/1=0$ would supply
an annihilator outside $\mathfrak m$. Hence every $a_i\in R$,
proving condition 2. Condition 2 gives $\Omega=R[T]$ and the
usual tensor isomorphisms, hence condition 3.

For the converse let $k=R/\mathfrak m=\mathbb F_q$ and form
\[
 S=R\oplus k\epsilon,\qquad
 (r,a\epsilon)(s,b\epsilon)
  =(rs,(\bar r b+\bar s a)\epsilon),\qquad \epsilon^2=0.
\]
The map $r\mapsto(r,0)$ is a unital inclusion. Writing a polynomial
uniquely as $F_0+\epsilon H$ gives
\[
 \Omega=R[T]\oplus\epsilon(T^q-T)k[T],
\]
because the null polynomials on $\mathbb F_q$ are exactly the
multiples of $T^q-T$; division by this monic polynomial proves
that assertion. In particular condition 2 fails. Put
$e=\epsilon(T^q-T)$. The module $ke$ is a direct $R$-summand of
$\Omega$: in the second summand divide by $T^q-T$ and extract
the constant coefficient. Hence $e\otimes e$ is nonzero in
$\Omega\otimes_R\Omega$, since
$ke\otimes_Rke\simeq k\otimes_Rk\simeq k$.
Its polynomial image is zero because $\epsilon^2=0$.
Condition 3 therefore fails as well. This proves both converses
and the complete equivalence. The finite $R=k=\mathbb F_2$
instance is checked explicitly in \texttt{check\_B\_domains.py}.
\end{proof}

\begin{saresult}{B.018}{Proved}{The free-module repair for every torsion-free extension of an infinite domain}
Let $R$ be an infinite domain and $R\subset S$ a unital inclusion
with $S$ torsion-free over $R$. If
$\Omega=\{F\in S[T]:F(R)\subset R\}$ is free over $R$, then
for every $n\ge1$ the natural map is an isomorphism
\[
 \Omega^{\otimes_R n}\xrightarrow{\sim}
 \{P\in S[T_1,\ldots,T_n]:P(R^n)\subset R\}.
\]
Consequently the two requested tensors and the compatible plethory
structure exist uniquely. This extends B.009 from a fraction field
to every torsion-free overring of an infinite domain; B.008 shows
why the word infinite cannot simply be omitted.
\end{saresult}
\begin{proof}
By B.014 replace the coefficient ring by
$L=S_{\mathrm{rat}}\subset K=\operatorname{Frac}(R)$.
The same reduction holds in any finite number of variables:
use a rectangular grid of distinct elements of $R$ and invert
the Vandermonde matrix in each variable. Every coefficient of an
integer-valued multivariable polynomial then lies in $K\cdot1$
and hence in $L$.

Let $\mathcal B$ be an $R$-basis of $\Omega$. Localization shows
that it is a $K$-basis of $K[T]$. It is in fact an $L$-basis of
$L[T]$: independence follows after embedding in $K[T]$, and
spanning follows because every monomial $T^n$ already lies in
$\Omega$ and thus has a finite expansion with coefficients in $R$.
Now expand $P(X,Y)\in L[X,Y]$, integer-valued on $R^2$, uniquely
as $\sum_i b_i(X)h_i(Y)$ using this $L$-basis. The coefficients
$h_i$ lie in $L[Y]$. At every $y\in R$, the polynomial $P(X,y)$
belongs to $\Omega$, so uniqueness in the $K$-basis forces
$h_i(y)\in R$. Thus $h_i\in\Omega$. This gives surjectivity.
The tensors of the $R$-basis remain linearly independent over $K$,
which proves injectivity. Iteration gives the assertion for every
$n$. The same transport of polynomial identities as in B.009
then gives the biring and plethory axioms, proving the conclusion.
\end{proof}

\begin{saresult}{B.019}{Proved}{A finite dual-basis criterion, including every projective case}
In the infinite-domain, torsion-free-extension setting of B.014,
assume the following condition on the $R$-module $\Omega$: for every
finite subset $E\subset\Omega$ there are finitely many
$b_i\in\Omega$ and $R$-linear maps $\lambda_i:\Omega\to R$ such that
\[
 x=\sum_i b_i\lambda_i(x)\qquad(x\in E).
\]
Then every multivariable tensor map in B.018 is an isomorphism,
and $\Omega$ has the unique compatible plethory structure.
In particular this holds if $\Omega$ is projective over $R$;
freeness is unnecessary.
\end{saresult}
\begin{proof}
Reduce to the coefficient overring $L\subset K$ by B.014 and
its grid argument. Every $\lambda_i$ extends uniquely to a
$K$-linear functional on $K\otimes_R\Omega=K[T]$.
It maps $L[T]$ into $L$: its value on each monomial is in $R$,
and linearity gives the assertion for arbitrary coefficients in $L$.

Let $P(X,Y)\in L[X,Y]$ be integer-valued on $R^2$. Choose a nonzero
$d\in R$ clearing the denominators of its finitely many coefficient
polynomials in $X$, so those coefficients of $dP$ belong to $R[X]$
inside $\Omega$. Choose the finite dual-basis identity on this finite
set. After extension to $K$, it holds on their $K$-linear span,
so
\[
 P(X,Y)=\sum_i b_i(X)h_i(Y),\qquad
 h_i(Y)=\lambda_{i,K}\bigl(P(X,Y)\bigr).
\]
The preceding observation gives $h_i\in L[Y]$. At $y\in R$,
$P(X,y)\in\Omega$, so $h_i(y)=\lambda_i(P(X,y))\in R$.
Thus $h_i\in\Omega$, proving surjectivity.

For injectivity, let $z=\sum_j x_j\otimes y_j$ have zero polynomial
image and choose a finite dual-basis identity on the $x_j$. Then
\[
 z=\sum_i b_i\otimes\left(\sum_j\lambda_i(x_j)y_j\right).
\]
Applying $\lambda_{i,K}$ in the first variable to the zero polynomial
image gives $\sum_j\lambda_i(x_j)y_j=0$ in $K[Y]$, hence in $\Omega$.
Thus $z=0$. The same argument one variable at a time proves both
assertions for all tensor powers. Polynomial identities then give
the plethory as in B.009.

Finally a projective module is a summand of a free module. Embed
$\Omega$ in a free module and project back. The finitely many free
coordinates occurring in a finite subset $E$ give coordinate
functionals $\lambda_i$ on $\Omega$ and projected free basis vectors
$b_i$ with exactly the required identity. This verifies the
hypothesis rather than assuming that an arbitrary torsion-free
module has a dual basis.
\end{proof}

\begin{saresult}{B.020}{Proved}{Every infinite Noetherian normal base with torsion-free coefficients}
Let $R$ be an infinite Noetherian normal domain and $R\subset S$
an arbitrary $R$-torsion-free unital extension. Then $\Omega$ is a
projective $R$-module, all the multivariable tensor maps are
isomorphisms, and the unique co-operations and compatible plethory
exist. More explicitly, put $K=\operatorname{Frac}(R)$ and
$L=S_{\mathrm{rat}}\subset K$. For $n\ge0$ put
\[
 w_q(n)=\sum_{a\ge1}\left\lfloor\frac{n}{q^a}\right\rfloor,
 \qquad
 I_n=\prod_{\substack{\operatorname{ht}\mathfrak p=1\\
                     |R/\mathfrak p|=q<\infty}}
                 \mathfrak p^{-w_q(n)}.
\]
Only finitely many factors are nontrivial. There are monic
$P_n\in R[T]$ of degree $n$ such that
\[
 \operatorname{Int}(R)=\bigoplus_{n\ge0}I_nP_n,
 \qquad
 \Omega=\bigoplus_{n\ge0}(I_n\cap L)P_n.
\]
Each fractional ideal $I_n\cap L$ is invertible. These formulas
include nonprincipal ideals and therefore do not assume a regular
free polynomial basis.
\end{saresult}
\begin{proof}
We first record precisely the commutative-algebra inputs. A
Noetherian normal domain is the intersection of its height-one
localizations in its fraction field, and each such localization
is a discrete valuation ring. These are the opened Stacks Project
Lemmas 10.157.6 and 10.119.7, tags 031T and 00PD
\cite{StacksNormalDomains}. A finitely presented module whose
localizations at all maximal ideals are free is finite projective,
by its Lemma 10.78.2, tag 00NX. We apply these results only with
exactly the stated Noetherian and normal hypotheses.

There are finitely many maximal ideals with residue field of
cardinality at most $n$. To prove this, for each prime power $q\le n$
let $J_q$ be the ideal generated by $x^q-x$, $x\in R$.
Every prime quotient of $R/J_q$ is a domain in which every element
is a root of $T^q-T$, and therefore is a finite field with at most
$q$ elements. All its primes are consequently maximal. A
Noetherian ring has finitely many minimal primes, so a
zero-dimensional Noetherian ring has finitely many primes.
Every maximal ideal with residue field $\mathbb F_q$ contains
$J_q$, proving finiteness. A height-one prime with finite residue
ring is maximal. Such a prime $\mathfrak p$ is invertible: at
$\mathfrak p$ it is principal in a discrete valuation ring, and at
every other maximal ideal its localization is the whole ring.
It is finitely presented because $R$ is Noetherian, so the cited
finite-projective criterion applies. In particular every $I_n$
in the statement is an invertible fractional ideal.

Here is the needed discrete-valuation calculation with proof.
Let $V$ be a discrete valuation ring with uniformizer $\pi$ and
residue field of cardinality $q$. Choose representatives
$c_0=0,c_1,\ldots,c_{q-1}$ of its residue classes. If
$n=\sum_j n_jq^j$, define $a_n=\sum_jc_{n_j}\pi^j$ and
$Q_n(T)=\prod_{i<n}(T-a_i)$. The first unequal digits show that
$a_i\equiv a_j\pmod{\pi^r}$ exactly when
$i\equiv j\pmod{q^r}$. Among $0,\ldots,n-1$, every residue class
modulo $q^r$ occurs at least $\lfloor n/q^r\rfloor$ times, and the
class of $n$ occurs exactly that often. Counting the factors
congruent to an arbitrary $x\in V$ modulo $\pi^r$ gives
\[
 v(Q_n(x))\ge w_q(n),\qquad v(Q_n(a_n))=w_q(n).
\]
Thus $Q_n/\pi^{w_q(n)}$ is integer-valued and has unit value at
$a_n$. Successively evaluating a polynomial at $a_0,a_1,\ldots$
shows that these polynomials are a $V$-basis of
$\operatorname{Int}(V)$: subtract its degree-zero coefficient,
then its coefficient of $Q_1/\pi^{w_q(1)}$, and so forth; at step
$n$ the diagonal value is a unit. The possible leading coefficients
in degree at most $n$ are exactly $\pi^{-w_q(n)}V$.

Apply this calculation at a height-one maximal ideal
$\mathfrak p$ of $R$ with finite residue field. A uniformizer in
$R_{\mathfrak p}$ can be chosen from $R$, since one generator of
the finitely generated ideal $\mathfrak p$ has valuation one.
The residue representatives can also be chosen from $R$.
Thus every node $a_i$ above belongs to $R$. The same successive
evaluation argument therefore bounds the leading coefficient of
any $F$ with $F(R)\subset R_{\mathfrak p}$ by
$-w_q(n)$; it does not assume a localization theorem for
integer-valued polynomials. At any other height-one prime,
$R/\mathfrak p$ is infinite, and Vandermonde evaluation at $n+1$
distinct residue classes makes every coefficient integral in
$R_{\mathfrak p}$. These observations show that the leading
coefficient of every degree-$n$ element of $\operatorname{Int}(R)$
belongs to $I_n$. Indeed these are exactly its height-one valuation
bounds: if $c$ satisfies them, then $cb\in R$ for every
$b\in I_n^{-1}$ by the intersection theorem, and
$(I_n^{-1})^{-1}=I_n$ since $I_n$ is invertible.

For the reverse inclusion, let $E_n$ be the finite set of primes
contributing nontrivially to $I_n$. For each $\mathfrak p\in E_n$
construct the monic $Q_{n,\mathfrak p}\in R[T]$ just described.
By the Chinese remainder theorem choose a monic $P_n\in R[T]$
of degree $n$ whose lower coefficients are congruent to those of
$Q_{n,\mathfrak p}$ modulo $\mathfrak p^{w_q(n)}$, simultaneously
for all $\mathfrak p\in E_n$. The relevant ideal powers are
pairwise comaximal. Its values have valuation at least $w_q(n)$
at each of those primes. Consequently $cP_n$ is integer-valued
for every $c\in I_n$: it is integral at every height-one prime,
and their intersection is $R$. This proves that $I_n$ is exactly
the leading-coefficient ideal. Subtracting $cP_n$ and descending
in degree proves
$\operatorname{Int}(R)=\bigoplus I_nP_n$; uniqueness follows from
the distinct monic degrees.

Because every $P_n$ is monic with coefficients in $R$, that same
descending-degree procedure shows that its expansion has all
coefficients in $L$ precisely when the original polynomial does.
This proves $\Omega=\bigoplus (I_n\cap L)P_n$.
Set $J_n=I_n\cap L$. It is finitely presented, being a submodule of
the finite fractional ideal $I_n$ over the Noetherian ring $R$.
Localization commutes with this intersection. At a maximal ideal
outside $E_n$, it gives $(J_n)_{\mathfrak m}=R_{\mathfrak m}$.
At $\mathfrak p\in E_n$, the ring $L_{\mathfrak p}$ is an overring
of the discrete valuation ring $R_{\mathfrak p}$ in $K$, and is
either that ring or $K$. To verify this dichotomy, an element of
negative valuation, multiplied by a unit and a suitable power of
$\pi$, supplies $\pi^{-1}$, and inverting $\pi$ gives $K$.
Accordingly $(J_n)_{\mathfrak p}$ is either $R_{\mathfrak p}$ or
$\pi^{-w_q(n)}R_{\mathfrak p}$. Every localization is therefore
free of rank one, so $J_n$ is invertible by the cited criterion.
The displayed direct sum is projective. B.019 now proves all the
tensor isomorphisms and the plethory assertion.
\end{proof}

\begin{saresult}{B.021}{Proved}{An explicit infinite-rank patching lemma}
Let $V\to D$ be a surjective ring homomorphism, let $k\to D$ be a
second homomorphism, and put $R=V\times_D k$. Fix a set $I$ and an
isomorphism $\alpha:D^{(I)}\to D^{(I)}$. Define
\[
 M_\alpha=\{(b,c)\in V^{(I)}\oplus k^{(I)}:
                  \bar b=\alpha(\bar c)\}.
\]
Then $M_\alpha$ is projective over $R$, even for infinite $I$.
More precisely $M_\alpha\oplus M_{\alpha^{-1}}$ is free over $R$.
\end{saresult}
\begin{proof}
All matrices below act on direct sums, so each column has finite
support. Matrix products are well defined: applying a matrix to a
finite-support vector again gives a finite-support vector.
In the endomorphism ring of $D^{(I)}$ one has
\[
 \begin{pmatrix}\alpha&0\\0&\alpha^{-1}\end{pmatrix}
 =
 \begin{pmatrix}1&\alpha\\0&1\end{pmatrix}
 \begin{pmatrix}1&0\\-\alpha^{-1}&1\end{pmatrix}
 \begin{pmatrix}1&\alpha\\0&1\end{pmatrix}
 \begin{pmatrix}0&-1\\1&0\end{pmatrix}.
\]
Direct block multiplication proves this equality. Lift the finitely
many entries in each column of $\alpha$ and $\alpha^{-1}$ arbitrarily
through $V\to D$, choosing zero lifts of zero entries. Replace the
two off-diagonal matrices in the displayed product by these lifts.
Each elementary block matrix is invertible over $V$, regardless of
whether the lifted off-diagonal matrices are inverses. Their product
is an automorphism $\Gamma$ of $V^{(I\sqcup I)}$ reducing to
$\operatorname{diag}(\alpha,\alpha^{-1})$.

The direct sum $M_\alpha\oplus M_{\alpha^{-1}}$ is precisely the
module patched with that block diagonal gluing map. Applying
$\Gamma^{-1}$ to its $V$-component and the identity to its
$k$-component identifies it with
\[
 \{(b,c)\in V^{(I\sqcup I)}\oplus k^{(I\sqcup I)}:
                       \bar b=\bar c\}=R^{(I\sqcup I)}.
\]
The last equality is coordinatewise; all supports are finite.
Thus $M_\alpha$ is a direct summand of a free module, as claimed.
No finite-rank assumption or infinite version of Nakayama's lemma
has entered the argument.
\end{proof}

\begin{saresult}{B.022}{Proved}{The complete equal-residue-field cusp family}
Let $k=\mathbb F_q$, let $m\ge1$, put
\[
 V=k[[t]],\qquad K=k((t)),\qquad R=k+t^mV.
\]
Then $R$ is a one-dimensional Noetherian local domain with fraction
field $K$, and $A=\operatorname{Int}(R)$ is a projective $R$-module.
Consequently all its multivariable tensor maps are isomorphisms
and it has the unique compatible plethory structure.
In particular $R=\mathbb F_2[[t^2,t^3]]$ satisfies these assertions.
The conjectured nonflatness for this example recorded in the opened
Elliott paper \cite[introduction, p.2]{ElliottPresentations2011}
is false: projectivity implies flatness.
\end{saresult}
\begin{proof}
The maximal ideal is $t^mV$, and the residue field is $k$;
an element with nonzero constant term has its inverse in $R$.
The ring is finite free over $k[[t^m]]$, with basis
$1,t^{m+1},\ldots,t^{2m-1}$ (for $m=1$, just $1$).
Thus it is Noetherian and one-dimensional, and it is a domain
inside $V$. Also $t^mV\subset R$ implies
$V\subset\operatorname{Frac}(R)$, hence its fraction field is $K$.
For $m=2$, the allowed exponents are exactly $0,2,3,4,\ldots$,
which gives the asserted power-series-ring identification.

First compute $B=\operatorname{Int}(R,V)\subset K[X]$ as a
$V$-module. List $k$ as $c_0=0,c_1,\ldots,c_{q-1}$. For the
base-$q$ expansion $n=\sum_{j\ge0}n_jq^j$ put
\[
 a_n=c_{n_0}+\sum_{j\ge1}c_{n_j}t^{m+j-1},\quad
 d_n=\prod_{i<n}(a_n-a_i),\quad
 b_n(X)=\frac{\prod_{i<n}(X-a_i)}{d_n}.
\]
All nodes belong to $R$ and are distinct. Congruence modulo
$t^r$ is determined by the first digit for $1\le r\le m$,
and by the first $r-m+1$ digits for $r>m$. Counting factors
congruent to a point $x\in R$ gives
\[
 v_t\!\left(\prod_{i<n}(x-a_i)\right)\ge
 m\left\lfloor\frac nq\right\rfloor+
 \sum_{j\ge2}\left\lfloor\frac n{q^j}\right\rfloor
 =v_t(d_n).
\]
The equality at $x=a_n$ follows because its residue class has
exactly the displayed floor counts. Hence every $b_n\in B$,
$b_n(a_n)=1$, and $b_n(a_i)=0$ for $i<n$. Successive Newton
interpolation at these nodes shows
$B=\bigoplus_{n\ge0}Vb_n$: at each step the next coefficient
is a value in $V$ and the diagonal value is exactly $1$.

Choose $Q=q^e\ge m$. For every $F\in B$, its values lie in $V$,
so $F^Q$ has values in $k+t^QV\subset R$ and belongs to $A$.
Moreover $F^Q-F$ has values in $tV$, hence belongs to $tB$.
Thus $A+tB=B$. Also
\[
 A\cap tB=t^mB:
\]
a value in $R\cap tV$ lies in $t^mV$, and division of the
polynomial by $t^m$ is still a polynomial in $K[X]$ with values
in $V$. Conversely $t^mB\subset A\cap tB$ is immediate.
It follows that $C=A/t^mB$ maps isomorphically, as a $k$-vector
space, onto $B/tB$.

Multiplication defines a map
\[
 \alpha:C\otimes_k(V/t^mV)\longrightarrow B/t^mB.
\]
Choose the $k$-basis of $C$ corresponding to the reductions of
the $b_n$. Relative to these bases, $\alpha$ is a column-finite
matrix congruent to the identity modulo $t$. Write it as $1+N$
with every entry of $N$ divisible by $t$. Then $N^m=0$, and
$1-N+N^2-\cdots+(-N)^{m-1}$ is an inverse. Thus $\alpha$ is an
isomorphism by a finite nilpotent calculation, not an appeal to
Nakayama for infinitely generated modules.

The ring $R$ is the fiber product $V\times_{V/t^mV}k$, and $A$
is exactly the module patched from the free $V$-module $B$ and
the free $k$-module $C$ by this $\alpha$. Indeed a polynomial in
$B$ has reduction in $C\subset B/t^mB$ precisely when its values
lie in $R$; this is also immediate from the definition of $C$ and
$t^mB\subset A$. B.021 makes $A$ projective. B.019, applied with
coefficient field $K$, gives all the tensor isomorphisms and the
plethory. Flatness follows directly because tensoring with a
direct summand of a free module preserves injections.
\end{proof}

\begin{saresult}{B.023}{Proved}{Finite residue extensions and arbitrary conductor gaps}
Let $k\subset l$ be finite fields of cardinalities $q$ and
$Q=q^f$, let $m\ge1$, and put
\[
 V=l[[t]],\qquad K=l((t)),\qquad R=k+t^mV.
\]
Then $R$ is a one-dimensional Noetherian local domain and
$\operatorname{Int}(R)$ is projective over $R$. All multivariable
tensor maps are isomorphisms, so its unique compatible plethory
exists. In particular this proves flatness also for
$\mathbb F_2+t\mathbb F_4[[t]]$, the second conjectured nonflat
example recorded in \cite[introduction, p.2]{ElliottPresentations2011}.
Neither example is the two-branch local domain of B.016.
\end{saresult}
\begin{proof}
Locality, the maximal ideal $t^mV$, and the fraction field $K$
are verified as in B.022. If $1,e_1,\ldots,e_{f-1}$ is a
$k$-basis of $l$, then $R$ is free of rank $mf$ over $k[[t^m]]$,
with basis
\[
 1;\quad e_jt^m\ (1\le j<f);\quad
 e_jt^{m+r}\ (0\le j<f,\ 1\le r<m),\qquad e_0=1.
\]
This follows by collecting the coefficients in each residue class
of exponents modulo $m$. Hence $R$ is Noetherian. For a nonzero
element $x\in R$, multiplication by $x$ on this finite free module
has nonzero determinant. Its characteristic equation puts that
determinant in $xR$. Thus any nonzero prime of $R$ meets
$k[[t^m]]$ nontrivially and lies over its maximal ideal; since $R$
is local and integral over this ring it has dimension one.

Put $B=\operatorname{Int}(R,V)$. The Newton construction of B.022
uses mixed digits here. List the elements of $k$ as $c_0,\ldots,c_{q-1}$
and those of $l$ as $d_0,\ldots,d_{Q-1}$, both starting with zero.
Write
\[
 n=n_0+q\sum_{j\ge1}n_jQ^{j-1},\qquad
 0\le n_0<q,\quad 0\le n_j<Q,
\]
and set
\[
 a_n=c_{n_0}+\sum_{j\ge1}d_{n_j}t^{m+j-1},\qquad
 b_n(X)=\frac{\prod_{i<n}(X-a_i)}{\prod_{i<n}(a_n-a_i)}.
\]
Exactly the same residue-class counting gives the denominator
valuation
\[
 m\left\lfloor\frac nq\right\rfloor+
 \sum_{j\ge1}\left\lfloor\frac n{qQ^j}\right\rfloor,
\]
and gives at least this valuation at every point of $R$.
The diagonal Newton values are $1$, so
$B=\bigoplus_{n\ge0}Vb_n$.

We carefully distinguish polynomials from their residue functions.
Evaluation gives an injection
\[
 D_0=B/tB\longrightarrow l^R:
\]
its kernel before quotienting consists exactly of polynomials with
values in $tV$, namely $tB$. Thus every $z\in D_0$ satisfies
$z^Q=z$. The map $\sigma:z\mapsto z^q$ is an automorphism of
this ring, is semilinear for the $q$-power automorphism of $l$,
and satisfies $\sigma^f=1$. Put $C_0=D_0^\sigma$.

Here is a direct finite-field descent proof:
\[
 C_0\otimes_k l\xrightarrow{\sim}D_0.
\]
Every finitely generated $l$-subalgebra of $D_0$ is a quotient of
$l[X_1,\ldots,X_r]/(X_i^Q-X_i)$, a finite product of copies of
$l$ by the Chinese remainder theorem. It is therefore itself a
finite product of copies of $l$. On such a product $\sigma$ is
coordinatewise $q$-power, since it is the actual power operation
on elements; it fixes each idempotent and its fixed ring is the
corresponding product of copies of $k$. Scalar extension gives
the asserted isomorphism on each finite subalgebra. These
subalgebras are stable under $\sigma$ and form a filtered union
of $D_0$. Taking that union proves the displayed descent statement.
In particular no unverified assertion about Frobenius acting on
polynomial coefficients has been substituted for its action on
residue polynomial functions.

Let $A=\operatorname{Int}(R)$. Its image in $D_0$ is exactly
$C_0$. Indeed its residue values lie in $k$, which is the fixed
field of $q$-power. Conversely, if $F\in B$ has residue function
in $C_0$, choose $N=Q^e\ge m$. Then $F^N$ has values in
$k+t^NV\subset R$ and has the same residue function as $F$.
The kernel of $A\to D_0$ is $A\cap tB=t^mB$, by the same
valuewise argument as in B.022. Hence
$C=A/t^mB\simeq C_0$ as $k$-algebras.

The multiplication map
\[
 C\otimes_k(V/t^mV)\longrightarrow B/t^mB
\]
reduces modulo $t$ to the proved descent isomorphism.
Choose bases of the free modules on both sides. Their reduced
matrix is a column-finite invertible matrix over $l$; lift it
and its inverse by constant coefficients in $V/t^mV$ to change
one basis. The resulting matrix is $1+N$ with all entries in
$t(V/t^mV)$, and $N^m=0$. Its finite geometric inverse proves
that the multiplication map itself is an isomorphism. This is
valid in infinite rank because all columns and all sums in each
matrix product remain finite.

Finally $R=V\times_{V/t^mV}k$. The polynomial module $A$ is the
pullback of the free $V$-module $B$ and the free $k$-module $C$
along the last isomorphism. B.021 makes it projective over $R$;
B.019 gives the tensor isomorphisms and the plethory. This proves
the whole family and both stated special cases.
\end{proof}

\begin{saresult}{B.024}{Proved}{Every finite-conductor subring of one discrete valuation ring}
Let $V$ be a discrete valuation ring with fraction field $K$,
uniformizer $\pi$, and finite residue field of cardinality $Q$.
Let $R\subset V$ be a unital subring containing $\pi^mV$ for some
$m\ge1$. Then $R$ is a one-dimensional Noetherian local domain,
$V$ is its finite normalization, and
$A=\operatorname{Int}(R)$ is projective over $R$. Consequently
all multivariable tensor maps are isomorphisms and $A$ has the
unique compatible plethory. This includes incomplete and mixed
characteristic rings, and arbitrary conductor gaps rather than
only the pure gaps of B.023.
\end{saresult}
\begin{proof}
Put $J=\pi^mV$ and $R_0=R/J\subset D=V/J$. The ring $D$ is finite.
Its finitely many representatives generate $V$ over $R$, since
$J\subset R$. If $x\in R$ has valuation zero, multiplication by
its residue is injective on the finite ring $R_0$, hence surjective.
There is $y\in R$ with $xy=1+j$, $j\in J$. The inverse of $1+j$
in $V$ lies in $1+J\subset R$, so $x$ is a unit of $R$.
Thus $R$ is local with maximal ideal $R\cap\pi V$.
Also $V\subset\operatorname{Frac}(R)$ because $\pi^mV\subset R$.

For any nonzero ideal $H\subset R$, choose $h\in H$ of least
valuation. Then $H\subset hV$ and $JhV\subset H$.
The quotient $hV/JhV$ is finite, while $JhV$ is a finite $R$-module
because $V$ is. Finitely many representatives therefore generate
$H$ over $R$. This proves that $R$ is Noetherian. Every nonzero
prime has finite domain quotient, since it contains some $JhV$;
that quotient is a field. The nonzero maximal ideal gives dimension
one. Finally $V$ is finite, hence integral, over $R$ and is
integrally closed in $K$, so it is exactly the normalization.

Set $B=\operatorname{Int}(R,V)$. It is free over $V$ by an explicit
greedy Newton construction. Choose $a_0\in R$. Given distinct
$a_0,\ldots,a_{n-1}$, choose $a_n\in R$ minimizing the finite value
\[
 v_\pi\left(\prod_{i<n}(a_n-a_i)\right).
\]
Such a minimum exists: $R$ is infinite, so some value is finite,
and finite valuations are nonnegative integers. Put
\[
 b_n(X)=\frac{\prod_{i<n}(X-a_i)}{\prod_{i<n}(a_n-a_i)}.
\]
Minimality gives $b_n(R)\subset V$, and $b_n(a_n)=1$.
Successive evaluation at the nodes proves
$B=\bigoplus_{n\ge0}Vb_n$, exactly as in B.022.

We next compute $B/JB$ completely. Let $E_N$ be the finite set of
congruence classes modulo $\pi^NV$ met by $R$, and let
$\mathcal F(R,D)$ be the union of the function rings $D^{E_N}$
under pullback along refinement. Evaluation induces
\[
 B/JB\xrightarrow{\sim}\mathcal F(R,D).
\]
Injectivity is immediate: a polynomial with values in $J$ is
$\pi^m$ times one in $B$. Every polynomial in $B$ is constant
modulo $J$ on a sufficiently fine congruence partition. Indeed,
if its coefficients have valuations at least $-h$, then
$v(F(x)-F(y))\ge v(x-y)-h$ for $x,y\in V$.

For surjectivity we construct all congruence-class indicators.
The polynomial
\[
 \delta(X)=\frac{X^Q-X}{\pi}
\]
is integer-valued on $V$. If $v(x-y)=r\ge1$, factor its numerator
difference to obtain
\[
 v(\delta(x)-\delta(y))=r-1:
\]
the other factor reduces to $Q\bar x^{Q-1}-1=-1$ in the residue
field. Hence the residues of
$X,\delta(X),\ldots,\delta^{\circ(N-1)}(X)$ distinguish precisely
the classes modulo $\pi^N$. For a chosen class, let $c_j\in V$
represent its $j$th residue digit and put
\[
 H(X)=\prod_{j=0}^{N-1}
       \left(1-(\delta^{\circ j}(X)-c_j)^{Q-1}\right).
\]
Modulo $\pi$, this is its characteristic function on $V$.
A sufficiently high $Q$-power makes it the characteristic function
modulo $\pi^m$. Explicitly $H^{Q^m}$ suffices: when $H\in\pi V$
its power has valuation at least $Q^m\ge m$; when $H=1+z$,
$v(z)\ge1$, each $Q$-power increases $v((1+z)^Q-1)$ by at least
one. For the latter assertion, factor $Q$ as a power of the
residue characteristic $p$ and use the binomial formula for
$(1+z)^p$: every intermediate coefficient is divisible by $p$,
and the final term has valuation $p\,v(z)\ge v(z)+1$.
The argument also covers characteristic $p$, where these
intermediate coefficients vanish. These indicators, multiplied
by arbitrary representatives in $V$ and added finitely, prove
surjectivity.

Put $C=A/JB$. The ideal $JB=\pi^mB$ is contained in $A$; no
identification of $JA$ with $JB$ is assumed.
Its image is exactly $\mathcal F(R,R_0)\subset\mathcal F(R,D)$.
One implication follows from values in $R$. Conversely a
polynomial whose residue function has values in $R_0$ has all
its actual values in $R$, since $J\subset R$; and the preceding
indicator construction lifts every such function. Thus
\[
 C=A/JB\simeq\mathop{\mathrm{colim}}_N R_0^{E_N}.
\]
This is a free $R_0$-module. To see it directly, select one child
in $E_{N+1}$ over each class in $E_N$. Restriction to the selected
children retracts the pullback
$R_0^{E_N}\to R_0^{E_{N+1}}$, with free kernel on the unselected
children. The union is the direct sum of the first finite free
module and these finite free kernels. Its rank is countably
infinite, since the numbers $|E_N|$ are unbounded for infinite $R$.

Tensor product commutes with these unions and with their finite
free terms. Consequently multiplication gives an isomorphism
\[
 C\otimes_{R_0}D\xrightarrow{\sim}B/JB.
\]
The ring square is $R=V\times_D R_0$, and $A$ is the module
pulled back from $B$ and $C$ along this isomorphism. Both modules
are free of countable rank, so B.021 applies and makes $A$
projective. B.019 proves the tensor and plethory conclusions.
In particular the proof did not assume that $JA=JB$ before the
patching argument, or use completeness or equal characteristic.
\end{proof}

\begin{saresult}{B.025}{Proved}{The branch criterion for finite normalization in dimension one}
Let $(R,\mathfrak m)$ be a one-dimensional Noetherian local domain
with finite residue field, fraction field $K$, and finite
normalization $T\subset K$. Then the following are equivalent:
\begin{enumerate}
\item $T$ is local;
\item every map $\operatorname{Int}(R)^{\otimes_R n}
\to\operatorname{Int}(R^n)$ is an isomorphism;
\item the polynomial identities of 153/004 give unique coaddition
and comultiplication tensors;
\item $F(XY)$ is a sum of terms $G(X)H(Y)$ with
$G,H\in\operatorname{Int}(R)$ for every $F\in\operatorname{Int}(R)$.
\end{enumerate}
If $T$ is not local, the proof gives an explicit higher-variable
polynomial obstruction, which rules out condition 4 by finite
iteration. Thus the criterion does not merely assert failure of
projectivity or of a sufficient condition.
\end{saresult}
\begin{proof}
The finite normalization is a one-dimensional Noetherian normal
semilocal domain. Its maximal ideals
$\mathfrak p_0,\ldots,\mathfrak p_s$ all contract to
$\mathfrak m$, and all their residue fields are finite.
Its localizations are discrete valuation rings, by the opened
normal-domain facts used in B.020. If $s=0$, it is a discrete
valuation ring. A nonzero conductor ideal exists: clear the
denominators of finitely many $R$-module generators of $T$.
In a discrete valuation ring it contains $\pi^mT$ for some $m$.
B.024 therefore proves condition 2. The implications
$2\Rightarrow3\Rightarrow4$ are immediate from B.009.

Suppose now $s\ge1$. We verify the conductor and approximation
steps explicitly. For each $i$, choose $\pi_i\in T$ congruent
to a valuation-one element modulo $\mathfrak p_i^2$, and
congruent to $1$ modulo every other maximal ideal. The Chinese
remainder theorem supplies it. Thus
$v_i(\pi_i)=1$ and $v_j(\pi_i)=0$ for $j\ne i$.
Localizing at all maximal ideals shows
$\mathfrak p_i=\pi_iT$; no general claim that a nonprincipal
ideal has a chosen generator is being used.
Write the nonzero conductor as
\[
 I=(R:T)=\prod_i\mathfrak p_i^{c_i}=aT,
 \qquad a=\prod_i\pi_i^{c_i}.
\]
The equality follows by localizing its two sides at the finitely
many discrete valuation rings. Every $c_i$ is positive. Otherwise
some conductor element lies outside $\mathfrak p_i$, hence outside
$\mathfrak m$ in $R$, so it is a unit of $R$. The conductor would
then contain $1$ and force $T=R$, contrary to $s\ge1$.
In particular $a\in I\subset R$ and $v_0(a)=c_0$.

Set $r=c_0+1$ and choose $N>r c_0$.
The unit $\pi_1$ has finite order $L$ in
$(T/\mathfrak p_0^N)^*$, since this quotient is finite.
For $n\ge1$ put
\[
 u_n=a\pi_1^{Ln}\in I\subset R.
\]
Then $u_n\equiv a\pmod{\mathfrak p_0^N}$ and
$v_1(u_n)\longrightarrow\infty$. For every finite family
$\mathcal H\subset\operatorname{Int}(R)$, sufficiently large $n$
gives
\[
 h(u_n)-h(0)\in\mathfrak p_1\cap R=\mathfrak m
 \qquad(h\in\mathcal H).
\]
Indeed every nonconstant term of each of the finitely many
polynomials tends to zero for $v_1$, and both
values lie in $R$.

Use the digit-indicator construction in B.024 at the discrete
valuation ring $T_{\mathfrak p_0}$ to produce a polynomial $H$
which is $1$ modulo $\mathfrak p_0$ on the class of $a^r$ modulo
$\mathfrak p_0^N$, and $0$ on the zero class. These classes are
distinct because $v_0(a^r)=r c_0<N$. In that construction the
residue representatives can be chosen in $T$, and its only
denominators are powers of $\pi_0$. Thus
$H\in T[\pi_0^{-1}][X]$, it is integer-valued in
$T_{\mathfrak p_0}$, and its coefficients are integral at every
other maximal ideal of $T$. Consequently $H(R)\subset T$.
Put
\[
 F=H-H(0),\qquad g=aF.
\]
Then $F(R)\subset T$, $F(0)=0$, and
$g(R)\subset aT=I\subset R$. Also
\[
 v_0(g(u_n^r))=c_0\quad\hbox{for every }n,
\]
since $u_n^r\equiv a^r\pmod{\mathfrak p_0^N}$ and $F$ has
unit residue on this class.

The polynomial $g(X_1\cdots X_r)$ is not in the image of the
$r$-fold tensor map. Otherwise write it as a finite sum of
products $\prod_{j=1}^r h_{ij}(X_j)$ and choose $n$ large for
all these factors in the displayed difference condition.
Its iterated difference on the $2^r$ points with coordinates
$0,u_n$ would be
\[
 \sum_i\prod_{j=1}^r(h_{ij}(u_n)-h_{ij}(0))\in\mathfrak m^r.
\]
The other expression gives exactly $g(u_n^r)$, because $g(0)=0$.
Every element of $\mathfrak m^r$ has $v_0\ge r>c_0$, a
contradiction. Finally, if condition 4 held, repeated binary
expansion, first separating $X_1$ from $X_2\cdots X_r$ and
then expanding every remaining factor, would put
$g(X_1\cdots X_r)$ in that image. The expansions are finite at
each of the $r-1$ steps, so no coherence or injectivity assumption
is needed for this implication. Thus condition 4 fails.
This proves all the equivalences and the negative direction.
\end{proof}

\begin{saresult}{B.026}{Proved}{All coefficient extensions for the one-branch local class}
Let $R\subset V$ satisfy B.024, with fraction field $K$, and let
$R\subset S$ be any $R$-torsion-free unital extension. Put
$L=S_{\mathrm{rat}}\subset K$ as in B.014. Then the identities
on 153/004 determine unique coaddition and comultiplication tensors
if and only if
\[
 L=R\quad\hbox{or}\quad L=K.
\]
In both cases the compatible plethory exists. A proper intermediate
coefficient overring $R\subsetneq L\subset K$ has a noninjective
tensor map. Thus even in this one-branch class the answer depends
on the specified coefficient ring, and does not grow monotonically
with that ring.
\end{saresult}
\begin{proof}
B.014 reduces the problem to $R\subset L\subset K$.
Any overring $L$ not contained in $V$ is $K$. Indeed, take
$b\in L$ of negative valuation and choose $N$ so large that
$c=b^{-N}\in\pi^mV\subset R$. Then $c^{-1}=b^N\in L$.
For every $x\in K$, a sufficiently high power of $c$ sends $x$
into $\pi^mV\subset R$. Hence $K=R[c^{-1}]\subset L$.
Every remaining overring is contained in the finite $R$-module
$V$, so it is finite over the Noetherian ring $R$.

Suppose $R\subsetneq L\subset V$. The conductor
$J=(R:L)$ contains $\pi^mV$, so $R/J$ is finite. Choose one
representative $r_i$ of each of its classes and put
$P(T)=\prod_i(T-r_i)$. It is monic and $P(R)\subset J$.
Therefore $P L[T]\subset\Omega(R,L)$, and monic division gives
\[
 \Omega(R,L)=E\oplus PL[T],\qquad
 E=\{H\in L[T]:\deg H<\deg P,\ H(R)\subset R\}.
\]
In particular $LP$ is a direct $R$-summand. Injectivity of
$\theta_2$ would force multiplication
$L\otimes_RL\to L$ to be injective, by restricting to this
summand, and hence an isomorphism.

A finite unital extension of a local ring with this property and
injective structure map is the identity. Here is a proof. After
base change to $k=R/\mathfrak m$, the finite-dimensional algebra
$C=L/\mathfrak mL$ still has injective multiplication
$C\otimes_kC\to C$. It is nonzero by Nakayama's elementary
finite-module argument. Choose a linear functional on $C$ taking
$1$ to $1$. Applying it to $a\otimes1=1\otimes a$ shows $C=k$.
Thus $L=R+\mathfrak mL$, so the finite module $L/R$ equals
$\mathfrak m(L/R)$ and is zero. For completeness, the finite-module
argument used twice says that if generators $x_i$ satisfy
$x_i=\sum a_{ij}x_j$ with $a_{ij}\in\mathfrak m$, the adjugate
matrix of $1-(a_{ij})$ annihilates the generators by a determinant
congruent to $1$ modulo $\mathfrak m$, hence a unit. This proves
$L=R$, a contradiction. Therefore the tensor map is not injective.

If $L=R$, the answer is the usual polynomial plethory.
If $L=K$, it is B.024. These cases exhaust all overrings and
prove the equivalence.

For a fully explicit noninjective example take
$R=\mathbb F_2+t^2\mathbb F_2[[t]]$ and $L=V=\mathbb F_2[[t]]$.
With $q(T)=T^2-T$ one has
\[
 \Omega(R,V)=R\oplus RT\oplus q(T)V[T].
\]
Division by $q$ proves this decomposition; the remainder has
coefficients in $R$ by evaluating at $0,1$. The tensor
\[
 z=(tq)\otimes q-q\otimes(tq)
\]
has zero polynomial image but is nonzero in the split summand
$Vq\otimes_RVq$. After tensoring this summand with $R/\mathfrak m$
it maps to $\tau-\sigma$ in
$\mathbb F_2[\tau,\sigma]/(\tau^2,\sigma^2)$, which is nonzero.
The attached exact program checks all $16$ elements of this finite
tensor algebra and its multiplication-to-the-diagonal map.
By contrast, with coefficient field $K$ the same base ring works
by B.022. This verifies the asserted dependence on coefficients.
\end{proof}


\begin{saresult}{B.027}{Proved}{Localization of every multivariable polynomial module in dimension one}
Let $R$ be a one-dimensional Noetherian domain, $K$ its fraction
field, and $R\subset L\subset K$ an arbitrary overring. For
$n\ge1$ put
\[
 \Omega_n(R,L)=\{F\in L[X_1,\ldots,X_n]:F(R^n)\subset R\}.
\]
For every maximal ideal $\mathfrak m$ one has, as submodules of
$K[X_1,\ldots,X_n]$,
\[
 \Omega_n(R,L)_{\mathfrak m}
   =\Omega_n(R_{\mathfrak m},L_{\mathfrak m}).
\]
In particular the tensor maps in B.009 localize to the corresponding
maps for $(R_{\mathfrak m},L_{\mathfrak m})$. No polynomial
localization property is assumed in this statement.
\end{saresult}
\begin{proof}
We first prove the required density of inputs. If $0\ne d\in R$,
then either $d$ is a unit in $R_{\mathfrak m}$ or
\[
 \sqrt{dR_{\mathfrak m}}=\mathfrak mR_{\mathfrak m}.
\]
Indeed every nonzero prime of the one-dimensional local domain
$R_{\mathfrak m}$ is its maximal ideal. Because that ideal is
finitely generated, some power
$\mathfrak m^NR_{\mathfrak m}$ is contained in
$dR_{\mathfrak m}$. Also
\[
 R\longrightarrow R_{\mathfrak m}/\mathfrak m^NR_{\mathfrak m}
\]
is surjective: the ring $R/\mathfrak m^N$ has the unique maximal
ideal $\mathfrak m/\mathfrak m^N$, so every denominator outside
$\mathfrak m$ is already invertible in this quotient.
It follows that every element of $R_{\mathfrak m}$ can be
approximated by an element of $R$ modulo $dR_{\mathfrak m}$.
The same assertion is immediate if $d$ is a unit.

Take $F\in\Omega_n(R,L)$ and choose $d\ne0$ with
$dF\in R[X_1,\ldots,X_n]$. Given
$x\in R_{\mathfrak m}^n$, choose $r\in R^n$ with
$x_j-r_j\in dR_{\mathfrak m}$ for each $j$.
The difference of the values of the polynomial $dF$ is a sum
of terms divisible by these coordinate differences. Thus
\[
 dF(x)-dF(r)\in dR_{\mathfrak m},
 \qquad F(x)-F(r)\in R_{\mathfrak m}.
\]
As $F(r)\in R$, we have $F(x)\in R_{\mathfrak m}$.
Dividing $F$ by any element outside $\mathfrak m$ preserves this
conclusion and puts its coefficients in $L_{\mathfrak m}$.
This proves the first inclusion.

Conversely take
$G\in\Omega_n(R_{\mathfrak m},L_{\mathfrak m})$.
Clear the finitely many coefficient denominators outside
$\mathfrak m$: there is $t\in R\setminus\mathfrak m$ with
$F=tG\in L[X_1,\ldots,X_n]$. Choose $d\ne0$ with
$dF\in R[X_1,\ldots,X_n]$ and form the ideal
\[
 H=(dF(r):r\in R^n)\subset R.
\]
It is finitely generated because $R$ is Noetherian, and
$H_{\mathfrak m}\subset dR_{\mathfrak m}$ because $F(r)$ lies
in $R_{\mathfrak m}$. For each of finitely many generators
$h_i$ choose $u_i\notin\mathfrak m$ with $u_i h_i\in dR$.
Their product $u$ satisfies $uH\subset dR$. Therefore
$uF(R^n)\subset R$, so $uF\in\Omega_n(R,L)$ and
\[
 G=\frac{uF}{ut}\in\Omega_n(R,L)_{\mathfrak m}.
\]
This proves the reverse inclusion with a uniform denominator,
not a separate denominator for each value. Localization commutes
with finite tensor products; all polynomial multiplication maps
are the evident ones inside $K[X_1,\ldots,X_n]$. This gives the
last assertion.
\end{proof}

\begin{saresult}{B.028}{Proved}{The global finite-normalization criterion}
Let $R$ be a one-dimensional Noetherian domain with fraction field
$K$ and finite normalization $T$ in $K$. The following are
equivalent:
\begin{enumerate}
\item for every maximal ideal $\mathfrak m$ with finite residue
field, $T_{R\setminus\mathfrak m}$ is local, equivalently exactly
one maximal ideal of $T$ lies above $\mathfrak m$;
\item all maps $\operatorname{Int}(R)^{\otimes_R n}
\to\operatorname{Int}(R^n)$ are isomorphisms;
\item the polynomial coaddition and comultiplication of 153/004
have unique tensor representatives;
\item $F(XY)$ has a finite expansion
$\sum_iG_i(X)H_i(Y)$ with $G_i,H_i\in\operatorname{Int}(R)$
for every $F\in\operatorname{Int}(R)$.
\end{enumerate}
When these hold, $\operatorname{Int}(R)$ has its unique compatible
plethory. Thus, for fraction-field coefficients, B.2(c) is completely
settled in dimension one with finite normalization. Only the
maximal ideals with finite residue field enter the criterion.
\end{saresult}
\begin{proof}
For any maximal ideal $\mathfrak m$, localization of the finite
normalization is the finite normalization of $R_{\mathfrak m}$.
Here is the integrality check. It is finite, hence integral, over
$R_{\mathfrak m}$ and is integrally closed as a localization of
$T$. Conversely if $x\in K$ is integral over $R_{\mathfrak m}$,
write a monic equation with its finitely many denominators outside
$\mathfrak m$. Choose their common multiple $s\notin\mathfrak m$.
Multiplying the equation for $x$ by $s^n$ gives a monic equation
for $sx$ over $R$, because each coefficient of $(sx)^i$ is the
original coefficient times $s^{n-i}$. Thus $sx\in T$ and
$x\in T_{R\setminus\mathfrak m}$. The maximal ideals of this
localized integral extension correspond exactly to those of $T$
above $\mathfrak m$.

Assume condition 1. If $R/\mathfrak m$ is infinite, the local
Vandermonde argument of B.017 gives
\[
 \operatorname{Int}(R_{\mathfrak m})=R_{\mathfrak m}[X],
 \qquad
 \operatorname{Int}(R_{\mathfrak m}^n)
       =R_{\mathfrak m}[X_1,\ldots,X_n].
\]
For the multivariable statement apply the same invertible
Vandermonde matrix successively in each variable. All local
tensor maps are therefore isomorphisms. If the residue field
is finite, condition 1 and B.025 give the same conclusion.
By B.027 the kernel and cokernel of each global tensor map
localize to zero at every maximal ideal, so they are zero.
Indeed a nonzero element of any module has a proper annihilator;
a maximal ideal containing it gives a localization in which
that element remains nonzero. This proves condition 2.
The implications $2\Rightarrow3\Rightarrow4$ follow from B.009.

Finally assume condition 4 and fix a maximal ideal
$\mathfrak m$ with finite residue field. B.027 writes any
$G\in\operatorname{Int}(R_{\mathfrak m})$ as $F/s$, with
$F\in\operatorname{Int}(R)$ and $s\notin\mathfrak m$.
An expansion of $F(XY)$, divided by $s$, gives an expansion of
$G(XY)$ over $R_{\mathfrak m}$. B.025 then forces its finite
normalization to be local. Thus condition 1 holds.
The plethory identities follow from the all-variable tensor
isomorphisms exactly as in B.009; no independent gluing of
operations or an unstated localization theorem is needed.
\end{proof}

\begin{saresult}{B.029}{Proved}{Coefficient extensions over the global one-branch class}
Let $R$ and $T$ satisfy B.028, and suppose the equivalent conditions
of that result hold. For any $R$-torsion-free unital extension
$R\subset S$, put $L=S_{\mathrm{rat}}\subset K$ as in B.014.
Then the following are equivalent:
\begin{enumerate}
\item $\Omega(R,S)$ has uniquely determined polynomial
coaddition and comultiplication tensors;
\item all its multivariable tensor maps are isomorphisms;
\item for every maximal ideal $\mathfrak m$ with finite residue
field, one has $L_{\mathfrak m}=R_{\mathfrak m}$ or
$L_{\mathfrak m}=K$.
\end{enumerate}
Whenever these hold the compatible plethory exists. The choices
$R_{\mathfrak m}$ and $K$ may differ at different maximal ideals.
\end{saresult}
\begin{proof}
B.014 identifies the polynomial module with $\Omega(R,L)$.
At a maximal ideal with infinite residue field it localizes,
by B.027 and the Vandermonde argument, to the ordinary polynomial
module over $R_{\mathfrak m}$; this works for any overring
$L_{\mathfrak m}$. At a maximal ideal with finite residue field,
the normalized local ring is a single discrete valuation ring
and its nonzero conductor contains a power of its maximal ideal.
Thus B.026 applies. Its proof gives two facts: the local tensor
maps are all isomorphisms for $L_{\mathfrak m}=R_{\mathfrak m}$
or $K$, and the second tensor map is noninjective in every other
case. A global injective tensor map stays injective on localization,
so condition 1 implies condition 3. Conversely condition 3 makes
all localized tensor maps isomorphisms; B.027 and the kernel and
cokernel argument in B.028 make them global isomorphisms.
The remaining implication and the plethory assertion are B.009.
\end{proof}


\begin{saresult}{B.030}{Proved}{The finite-extension obstruction with a cofinite conductor}
Let $R\subset S$ be an injective unital extension of nonzero
commutative rings, with $R$ local and $S$ finite as an $R$-module.
Suppose its conductor
\[
 J=\{r\in R:rS\subset R\}
\]
has finite quotient $R/J$. Then the second polynomial tensor map
for $\Omega(R,S)$ is injective if and only if $S=R$.
Consequently unique polynomial coaddition and comultiplication
occur exactly in that case.
\end{saresult}
\begin{proof}
Choose one representative of each class of $R/J$ and let $P$ be
the monic product of the corresponding linear factors. Then
$P(R)\subset J$, so $PS[X]\subset\Omega(R,S)$.
Division by the monic polynomial $P$ shows
\[
 \Omega(R,S)=E\oplus PS[X],
 \quad E=\{F\in\Omega(R,S):\deg F<\deg P\}.
\]
Thus $SP$ is a direct $R$-summand. Restriction of the tensor map
to $SP\otimes_RSP$ is multiplication
$S\otimes_RS\to S$, followed by multiplication by the monic
polynomial $P(X)P(Y)$. Injectivity of the tensor map would make
$S\otimes_RS\to S$ injective and hence an isomorphism.

Let $k$ be the residue field of $R$. The finite-dimensional
$k$-algebra $C=S\otimes_Rk$ is nonzero by the finite-module
Nakayama argument given explicitly in B.026. Its multiplication
would be an isomorphism $C\otimes_kC\to C$, so
$(\dim_k C)^2=\dim_k C$ and $C=k$.
It follows that the finite module $S/R$ satisfies
$S/R=\mathfrak m(S/R)$ and is zero by the same argument.
Thus $S=R$. Conversely for $S=R$ the module is $R[X]$ and all
polynomial tensor maps are isomorphisms.
\end{proof}

\begin{saresult}{B.031}{Proved}{All local branches and all torsion-free coefficient extensions}
Let $(R,\mathfrak m)$ be a one-dimensional Noetherian local domain
with finite residue field, fraction field $K$, and finite
normalization $T$. Let $R\subset S$ be any $R$-torsion-free
unital extension, and put $L=S_{\mathrm{rat}}$.
Then the polynomial coaddition and comultiplication tensors for
$\Omega(R,S)$ are uniquely determined if and only if
\[
 L=R\quad\hbox{or}\quad (T\hbox{ is local and }L=K).
\]
In these cases every multivariable tensor map is an isomorphism
and the unique compatible plethory exists. If $T$ has more than
one maximal ideal and $L\not\subset T$, even the existence of
multiplicative tensors fails.
\end{saresult}
\begin{proof}
Reduce to $\Omega(R,L)$ by B.014. If $T$ is local this is B.026,
using the nonzero conductor furnished by finite normalization.
Suppose $T$ has at least two maximal ideals. The notation and
conductor computation of B.025 apply:
\[
 \mathfrak p_i=\pi_iT,\qquad
 I=(R:T)=aT,\qquad v_i(a)=c_i>0.
\]
Every $T/\mathfrak p_i^{c_i}$ is finite, so $T/I$ and $R/I$
are finite. If $R\subsetneq L\subset T$, then $L$ is a finite
$R$-module because $R$ is Noetherian and $T$ is finite.
Its conductor contains $I$, and B.030 makes its second tensor
map noninjective. This handles every proper intermediate
coefficient ring contained in the normalization.

Now suppose $L\not\subset T$ and choose $b\in L\setminus T$.
At least one valuation $v_i(b)$ is negative. Define
\[
 c=\prod_{v_j(b)<0}\pi_j^{-v_j(b)},\qquad d=bc.
\]
Both $c$ and $d$ lie in $T$: all their valuations are nonnegative,
and $T$ is the intersection of its localizations inside $K$.
The latter assertion also follows directly from the ideal of
denominators: if $x\notin T$, the proper ideal
$\{t\in T:tx\in T\}$ is contained in a maximal ideal whose
localization still excludes $x$. At each maximal ideal at least
one of $c,d$ is a unit. Hence $(c,d)=T$, and there are $p,q\in T$
with $pc+qd=1$. Therefore
\[
 \frac1c=p+qb\in T[b],\qquad
 \pi_i^{-1}=(c/\pi_i)c^{-1}\in T[b]
       \quad\hbox{if }v_i(b)<0.
\]
This proves the needed inverse inclusion; negative valuation
alone has not been mistaken for a globally invertible element.

Relabel such an index as $0$ and choose another branch as $1$.
The construction in B.025 produces
\[
 H\in T[\pi_0^{-1}][X],\qquad
 g=a\bigl(H-H(0)\bigr)\in\operatorname{Int}(R),
\]
for which $g(X_1\cdots X_{c_0+1})$ has no tensor expansion even
with factors from the full module $\operatorname{Int}(R)$.
Every coefficient of $g$ lies in $aT[b]\subset R[b]\subset L$,
because $aT\subset R$. Thus $g\in\Omega(R,L)$.
If all multiplicative tensors existed for this smaller module,
repeated binary expansion would give the forbidden finite
$(c_0+1)$-variable expansion, exactly as in B.025.
They therefore do not all exist.

The only remaining possibility for a multibranch normalization
is $L=R$, for which the module is $R[X]$. Combined with the
one-branch case this proves the assertion in full.
\end{proof}

\begin{saresult}{B.032}{Proved}{Complete coefficient classification in dimension one with finite normalization}
Let $R$ be a one-dimensional Noetherian domain, let $K$ be its
fraction field, and suppose its normalization $T$ in $K$ is a
finite $R$-module. For any $R$-torsion-free unital extension
$R\subset S$, put $L=S_{\mathrm{rat}}\subset K$.
The following are equivalent:
\begin{enumerate}
\item the polynomial identities on 153/004 determine unique
coaddition and comultiplication tensors for $\Omega(R,S)$;
\item every multivariable tensor map for $\Omega(R,S)$ is an
isomorphism onto the corresponding integer-valued polynomials;
\item for every maximal ideal $\mathfrak m$ with finite residue
field, either
\[
 L_{\mathfrak m}=R_{\mathfrak m},
 \quad\hbox{or}\quad
 L_{\mathfrak m}=K\ \hbox{ and }\ T_{R\setminus\mathfrak m}
       \hbox{ is local}.
\]
\end{enumerate}
These conditions also give the unique compatible plethory.
There is no restriction at a maximal ideal with infinite residue
field. The criterion includes all coefficient overrings and all
normalization branch counts in the stated class.
\end{saresult}
\begin{proof}
The rational-coefficient reduction of B.014 works in every finite
number of variables by applying its Vandermonde matrices
successively in each variable. Thus all modules and tensor maps
can be computed with $L$ in place of $S$.
B.027 identifies their localizations at every maximal ideal.
At a maximal ideal with infinite residue field the Vandermonde
argument of B.017 gives the usual polynomial module and tensor
isomorphisms. At a maximal ideal with finite residue field,
B.031 gives exactly the alternatives in condition 3.

Both injectivity and the existence of the two substituted
polynomial tensors pass from the global module to its localization:
write a local polynomial as $F/s$ by B.027 and divide a tensor
expansion of $F(X+Y)$ or $F(XY)$ by $s$.
Therefore condition 1 forces all the local alternatives and
implies condition 3. Conversely condition 3 makes every local
tensor map an isomorphism. Its global kernel and cokernel vanish
by the localization test proved in B.028, so condition 2 follows.
B.009 gives condition 1 and all the plethory identities from
condition 2. This proves the full equivalence.
\end{proof}


\begin{saresult}{B.033}{Proved}{Polynomial localization in every Noetherian dimension}
The conclusion of B.027 holds for every Noetherian domain $R$,
without a dimension restriction, and for localization at any
multiplicatively closed subset $U\subset R\setminus\{0\}$:
\[
 U^{-1}\Omega_n(R,L)=\Omega_n(U^{-1}R,U^{-1}L).
\]
The inclusion from left to right holds for every domain, without
Noetherianity.
\end{saresult}
\begin{proof}
We prove the inclusion by induction on the total degree $d$ of
$F\in\Omega_n(R,L)$. A constant such polynomial belongs to $R$.
For $s\in U$ set
\[
 G_s(X_1,\ldots,X_n)
   =s^dF(X_1,\ldots,X_n)-F(sX_1,\ldots,sX_n).
\]
All terms of total degree $d$ cancel; hence $G_s$ has smaller
degree and still takes $R^n$ to $R$. Every tuple in $(U^{-1}R)^n$
can be written $a/s$ with a common $s\in U$ and $a\in R^n$.
The induction hypothesis and
\[
 s^dF(a/s)=G_s(a/s)+F(a)
\]
show that $F(a/s)\in U^{-1}R$. Dividing $F$ by an additional
element of $U$ preserves this property. This proves the inclusion
without any approximation or dimension assumption.

For the converse, take a polynomial with coefficients in
$U^{-1}L$ and values in $U^{-1}R$. Clear its finitely many
coefficient denominators by $t\in U$ to obtain $F\in L[X]$,
where $X$ denotes all $n$ variables. Choose $d\ne0$ with
$dF\in R[X]$. The ideal $(dF(r):r\in R^n)$ is finitely generated
and is contained in $dU^{-1}R$ after localization.
Multiplying by one element $u\in U$ puts every generator in $dR$,
so $uF\in\Omega_n(R,L)$. This is exactly the uniform-denominator
argument of B.027, and proves the converse.
\end{proof}

\begin{saresult}{B.034}{Proved}{The punctured-spectrum coefficient ring and the depth test}
Let $(R,\mathfrak m)$ be a nonfield Noetherian local domain with
fraction field $K$. Define the explicitly specified overring
\[
 R^\sharp=\{x\in K:\mathfrak m^Nx\subset R
                    \text{ for some }N\ge0\}.
\]
Then
\[
 R^\sharp=\bigcap_{\mathfrak p\ne\mathfrak m}R_{\mathfrak p}
 \quad\hbox{inside }K,
 \qquad
 R^\sharp=R\ \Longleftrightarrow\ \operatorname{depth}R\ge2.
\]
If the residue field is finite, and $R\subset L\subset K$ is an
overring, put $E=L\cap R^\sharp$. Then
\[
 \Omega_n(R,L)=\Omega_n(R,E)\quad(n\ge1),
 \qquad
 \Omega_1(R,L)=R[X]\ \Longleftrightarrow\ E=R.
\]
Thus the ring controlling whether nonpolynomial operations can
occur is determined without referring to tensor maps.
\end{saresult}
\begin{proof}
The displayed set is closed under addition and multiplication,
using a common sufficiently high power of $\mathfrak m$.
For $x\in K$ its nonzero denominator ideal
$I_x=\{r\in R:rx\in R\}$ detects membership in every
$R_{\mathfrak p}$. If $x$ belongs to all localizations with
$\mathfrak p\ne\mathfrak m$, either $I_x=R$ or
$\sqrt{I_x}=\mathfrak m$. Noetherianity then gives
$\mathfrak m^N\subset I_x$. Conversely if
$\mathfrak m^Nx\subset R$ and $\mathfrak p\ne\mathfrak m$,
choose $a\in\mathfrak m\setminus\mathfrak p$; the equality
$x=(a^Nx)/a^N$ places $x$ in $R_{\mathfrak p}$.
This proves the intersection formula.

Suppose $a,b\in\mathfrak m$ is a regular sequence. If
$R^\sharp\ne R$, choose $x\notin R$ with
$\mathfrak m x\subset R$: starting with any element of
$R^\sharp\setminus R$, use a least annihilating power and a
suitable element of the preceding power. Then $ax,bx\in R$ and
$b(ax)=a(bx)$. The nonzerodivisor $b$ on $R/aR$ forces
$ax\in aR$; cancellation by nonzero $a$ gives $x\in R$,
a contradiction. Hence depth at least two implies $R^\sharp=R$.
Conversely, if no regular sequence of length two exists, choose
$0\ne a\in\mathfrak m$. Every element of $\mathfrak m$ is a
zerodivisor on $R/aR$. The zerodivisors of this finite Noetherian
module form the union of its finitely many associated primes.
Prime avoidance forces $\mathfrak m$ to be one of those primes.
Thus there is $z\notin aR$ with $\mathfrak m z\subset aR$;
$x=z/a$ belongs to $R^\sharp\setminus R$.
The associated-prime facts used in this short argument are the
opened Stacks Project, Section 10.63, Lemmas 10.63.5 and 10.63.9,
and the depth-zero formulation in Lemma 10.72.5.

For the coefficient assertion take $F\in\Omega_n(R,L)$ and a
prime $\mathfrak p\ne\mathfrak m$. The domain $R/\mathfrak p$
is not a field and is therefore infinite. Choose as many elements
of $R$ with distinct residues modulo $\mathfrak p$ as the degree
in each variable requires. The corresponding Vandermonde
matrices are invertible over $R_{\mathfrak p}$. Applying them
to the values of $F$ proves that every coefficient lies in
$R_{\mathfrak p}$. The intersection formula puts every coefficient
in $E$, proving the module equality. If $E=R$ this is the usual
polynomial ring.

Conversely let $x\in E\setminus R$. Choose $N$ with
$\mathfrak m^Nx\subset R$. The quotient $R/\mathfrak m^N$ is
finite: its finite filtration has finite-dimensional vector-space
quotients over the finite residue field. The monic polynomial
$P$ formed from all of its residue representatives satisfies
$P(R)\subset\mathfrak m^N$. Hence $xP\in\Omega_1(R,L)$, but
its leading coefficient $x$ is not in $R$. This proves the last
equivalence.
\end{proof}

\begin{saresult}{B.035}{Proved}{The higher-dimensional local classification with finite normalization}
Let $(R,\mathfrak m)$ be a Noetherian local domain of dimension
at least two with finite residue field, fraction field $K$, and
finite normalization $T$. Let $R\subset L\subset K$ and set
$E=L\cap R^\sharp$ as in B.034. The following are equivalent:
\begin{enumerate}
\item the polynomial coaddition and comultiplication tensors
for $\Omega(R,L)$ are uniquely determined;
\item $\Omega(R,L)=R[X]$;
\item $E=R$.
\end{enumerate}
When $E\ne R$ and $E\subset T$, the second tensor map fails
injectivity. When $E\not\subset T$, multiplicative tensors fail
to exist for all elements. In particular, for $L=K$ the answer
is exactly $\operatorname{depth}R\ge2$.
\end{saresult}
\begin{proof}
B.034 identifies the module with $\Omega(R,E)$ and proves
$2\Longleftrightarrow3$; ordinary polynomials prove
$2\Rightarrow1$. Suppose first $R\subsetneq E\subset T$.
As $R$ is Noetherian and $T$ is finite, $E$ is a finite
$R$-module. Choose finitely many module generators in $R^\sharp$;
a common power $\mathfrak m^N$ sends all of them into $R$.
Thus the conductor of $R\subset E$ contains $\mathfrak m^N$
and has finite quotient. B.030 makes its second tensor map
noninjective.

It remains to treat $E\not\subset T$. The finite normalization
$T$ is a Noetherian normal semilocal domain. Its maximal ideals
are finite in number and all lie over $\mathfrak m$.
Let $\mathcal D$ be the set of its maximal ideals of height one.
For $\mathfrak p\in\mathcal D$ the ring $T_{\mathfrak p}$ is a
discrete valuation ring. Choose $\pi_{\mathfrak p}$ congruent to
a valuation-one element modulo $\mathfrak p^2$, and congruent to
$1$ at every other maximal ideal. Localizing shows
\[
 \mathfrak p=\pi_{\mathfrak p}T;
\]
this uses a principal ideal only at these height-one maximal
ideals, not at the other branches.

Take $b\in E\setminus T$. Since a Noetherian normal domain is
the intersection of its height-one valuation rings, some
height-one prime has negative valuation on $b$. Any such prime
must be maximal. Indeed, for a nonmaximal height-one prime
$\mathfrak q$ of $T$, its contraction in $R$ differs from
$\mathfrak m$ by integrality, and B.034 gives
$b\in R_{\mathfrak q\cap R}\subset T_{\mathfrak q}$.
Therefore all negative valuations of $b$ are among
$\mathcal D$. Put
\[
 c=\prod_{\substack{\mathfrak p\in\mathcal D\\v_{\mathfrak p}(b)<0}}
          \pi_{\mathfrak p}^{-v_{\mathfrak p}(b)},\qquad d=bc.
\]
Then $c,d\in T$ by the height-one intersection formula.
At a maximal ideal occurring in the product, $d$ is a unit;
at every other maximal ideal, $c$ is a unit. Thus $(c,d)=T$,
and the Bezout argument of B.031 proves
\[
 \pi_{\mathfrak p}^{-1}\in T[b]
       \quad\text{whenever }v_{\mathfrak p}(b)<0.
\]
Fix one such branch $\mathfrak p_0$ and abbreviate its uniformizer
as $\pi$. There is another maximal ideal $\mathfrak p_1$ of
$T$: otherwise $T$ would be a one-dimensional local domain,
contrary to $\dim T=\dim R\ge2$ for an integral extension.
No dimension assumption on $T_{\mathfrak p_1}$ is made.

The nonzero conductor $I=(R:T)$ is contained in $\mathfrak m$,
since otherwise $T=R$ would be local. Choose any nonzero $a\in I$
and write $c_0=v_0(a)>0$, $r=c_0+1$, and $N>r c_0$.
The digit construction of B.024 supplies
\[
 H\in T[\pi^{-1}][X]
\]
whose residue at $\mathfrak p_0$ is $1$ on the class of $a^r$
modulo $\mathfrak p_0^N$, and $0$ on the zero class.
It is integer-valued at $T_{\mathfrak p_0}$; at every other
maximal ideal its coefficients are integral because $\pi$ is
a unit. Thus $H(R)\subset T$. Set
\[
 g=a(H-H(0)).
\]
Its values lie in $I\subset R$, and its coefficients lie in
$aT[b]\subset R[b]\subset L$. Hence $g\in\Omega(R,L)$.

We prove the finite-difference condition needed for the
obstruction, allowing the second branch to have higher dimension.
Given any finite family $\mathcal H\subset\operatorname{Int}(R)$,
choose a common nonzero denominator $e\in R$ for all its
polynomial coefficients. Put
\[
 e'=e/\pi^{v_0(e)}\in T.
\]
The membership follows from the height-one intersection formula,
and $e'\notin\mathfrak p_0$.
The ideals $e'\mathfrak p_1$ and $\mathfrak p_0^N$ are comaximal,
so there is $x\in e'\mathfrak p_1$ with
$x\equiv1\pmod{\mathfrak p_0^N}$. Let $u=ax\in I\subset R$.
Then $u\equiv a\pmod{\mathfrak p_0^N}$, while
\[
 u/e\in\mathfrak p_1T_{\mathfrak p_1}.
\]
The latter assertion uses that $\pi$ is a unit at
$\mathfrak p_1$ and $x/e'\in\mathfrak p_1$.
Each nonconstant term of $h(u)-h(0)$ is therefore in
$\mathfrak p_1T_{\mathfrak p_1}$, and its total value lies in
$R$. Consequently
\[
 h(u)-h(0)\in\mathfrak p_1\cap R=\mathfrak m
        \quad(h\in\mathcal H).
\]
At the chosen valuation branch,
$v_0(g(u^r))=c_0<r$.
If $g(X_1\cdots X_r)$ had a finite tensor expansion, apply this
construction to all its factors. Its iterated difference on
$\{0,u\}^r$ would lie in $\mathfrak m^r$ and hence have
valuation at least $r$, a contradiction. Repeated binary
expansion again shows that multiplicative tensors cannot exist
for every element of $\Omega(R,L)$.
This excludes the final case and proves the theorem.
\end{proof}

\begin{saresult}{B.036}{Proved}{Complete Noetherian finite-normalization classification in every dimension}
Let $R$ be a nonfield Noetherian domain whose normalization $T$
in its fraction field $K$ is finite. Let $R\subset S$ be any
$R$-torsion-free unital extension and put $L=S_{\mathrm{rat}}$.
For a maximal ideal $\mathfrak m$, write $R_m=R_{\mathfrak m}$,
$L_m=L_{\mathfrak m}$ and $T_m=T_{R\setminus\mathfrak m}$.
The polynomial identities of 153/004 give unique coaddition and
comultiplication tensors if and only if the following holds
at every maximal ideal with finite residue field:
\begin{enumerate}
\item if $\dim R_m=1$, then $L_m=R_m$, or $L_m=K$ and $T_m$
is local;
\item if $\dim R_m\ge2$, then
\[
 L_m\cap R_m^\sharp=R_m,
 \qquad
 R_m^\sharp=\{x\in K:(\mathfrak mR_m)^Nx\subset R_m
                         \text{ for some }N\ge0\}.
\]
\end{enumerate}
In every positive case all multivariable tensor maps are
isomorphisms and the unique compatible plethory exists.
For fraction-field coefficients $S=K$, this criterion is:
\begin{enumerate}
\item a single normalization branch at each finite-residue
maximal ideal of height one;
\item depth at least two at every finite-residue maximal ideal
of height at least two.
\end{enumerate}
No hypothesis or restriction is imposed at maximal ideals with
infinite residue field.
\end{saresult}
\begin{proof}
Use B.014, in all variables, to replace $S$ by $L$.
B.033 identifies every localized polynomial module and tensor
map. At infinite-residue maximal ideals use the Vandermonde
argument. At finite-residue maximal ideals of height one use
B.031; at those of larger height use B.035.
These local results prove necessity because injectivity and
existence of the substituted tensors localize, as in B.032.
Conversely their positive cases give all local tensor
isomorphisms. The global kernels and cokernels vanish by the
localization test in B.028, yielding all global isomorphisms.
B.009 gives the unique plethory.
Finally B.034 identifies $R_m^\sharp=R_m$ with depth at least
two. This proves the stated fraction-field specialization.
\end{proof}


\begin{saresult}{B.037}{Proved}{An explicit two-dimensional tensor kernel}
Put
\[
 T=\mathbb F_2[x,y]_{(x,y)},\qquad
 I=(x,y)^2T,\qquad R=\mathbb F_2+I,\qquad K=\mathbb F_2(x,y).
\]
Then $R$ is a two-dimensional Noetherian local domain with residue
field $\mathbb F_2$ and finite normalization $T$. With $q(X)=X^2-X$,
\[
 \operatorname{Int}(R)=R\oplus RX\oplus q(X)T[X].
\]
The nonzero tensor
\[
 z=(xq)\otimes q-q\otimes(xq)
\]
has zero polynomial image. Thus even fraction-field coefficients
do not give uniquely determined polynomial tensors in this
explicit higher-dimensional case, essentially of finite type over
$\mathbb F_2$.
\end{saresult}
\begin{proof}
Let $A=\mathbb F_2[x^2,y^2]_{(x^2,y^2)}$. Then
$T=A\oplus Ax\oplus Ay\oplus Axy$: in characteristic two the
square of every denominator with nonzero constant term belongs
to $A^*$, so this equality holds after the stated localizations.
The ring $R$ is generated as an $A$-module by
\[
 1,\ xy,\ x^3,\ xy^2,\ x^2y,\ y^3.
\]
Indeed an $Ax$ or $Ay$ coefficient must have zero constant term,
and the maximal ideal of $A$ is $(x^2,y^2)A$.
Thus $R$ is a finite $A$-module, and every ideal of $R$ is
finitely generated already as an $A$-module; this proves
Noetherianity. Also $T$ is finite over $R$, and both have
fraction field $K$, since $x=x^3/x^2$ and $y=y^3/y^2$.
Dimension is two by the integral dimension theorem.
The quotient $R/I$ is $\mathbb F_2$, and any element of $R$
outside $I$ has inverse in $1+I$ up to its constant unit;
thus $R$ is local. The normal ring $T$, a localization of a
polynomial unique factorization domain, is its normalization.

Every height-one prime of $T$ is nonmaximal and contracts to a
nonmaximal prime of $R$. Vandermonde interpolation at distinct
residues of $R$ puts every coefficient of an integer-valued
polynomial in the corresponding localization of $T$.
Intersecting these valuation rings gives coefficients in $T$.
Thus $\operatorname{Int}(R)=\Omega(R,T)$.
For $r=a+i$ with $a\in\mathbb F_2$ and $i\in I$, one has
$q(r)=i^2-i\in I$, so $qT[X]\subset\operatorname{Int}(R)$.
Monic division by $q$ leaves a linear remainder; evaluation at
$0,1$ forces both its coefficients into $R$. This proves the
displayed direct sum.

The summand $Tq$ is split over $R$, so $Tq\otimes_RTq$ injects
as a summand of the full tensor module. After base change to
$R/I=\mathbb F_2$ it becomes
\[
 D\otimes_{\mathbb F_2}D,
 \qquad D=\mathbb F_2[x,y]/(x,y)^2.
\]
The image of $z$ is $x\otimes1-1\otimes x$, a nonzero difference
of two of its nine basis vectors. Its polynomial image is zero
because both terms equal $xq(X)q(Y)$.
The attached exact program checks all $512$ elements and
$262144$ multiplication pairs of this finite tensor algebra;
the diagonal map has image of size $8$ and kernel of size $64$.
It also checks the null-polynomial/conductor identity on all
$256$ possible $R$-inputs modulo $(x,y)^4$ and all $1024$
coefficient classes in that quotient.
\end{proof}


\begin{saresult}{B.038}{Proved}{A derivation obstruction to tensor injectivity}
Let $R$ be a domain with fraction field $K$, let $V\subset K$ be
a discrete valuation overring of $R$, and let $D:K\to K$ be a
derivation satisfying $D(R)\subset V$. Suppose there are
$u_n\in R$ with $v(u_n)\to\infty$ and $D(u_n)=1$.
For $A=\operatorname{Int}(R)$ the map
\[
 \lambda:A\longrightarrow V,\qquad F\longmapsto F'(0)
\]
is well-defined. There is an $R$-balanced bilinear map
\[
 \beta:A\times A\longrightarrow K/V,
 \qquad \beta(F,G)=D(F'(0))G'(0)+V.
\]
If $P\in R[X]$, $c\in K$, $cP\in A$ and
$D(c)P'(0)^2\notin V$, then
\[
 (cP)\otimes P-P\otimes(cP)
\]
is a nonzero tensor with zero polynomial image.
\end{saresult}
\begin{proof}
Apply $D$ to $F(u_n)\in R$. The polynomial chain rule gives
\[
 D(F(u_n))=(DF)(u_n)+F'(u_n),
\]
where $DF$ means applying $D$ to the coefficients.
Its constant coefficient is $D(F(0))\in V$.
All nonconstant terms of $(DF)(u_n)$ have nonnegative valuation
for sufficiently large $n$, since their coefficients are fixed
and $v(u_n)\to\infty$. Hence $(DF)(u_n)\in V$ eventually.
The left side lies in $V$, so $F'(u_n)\in V$.
The nonconstant terms of $F'(u_n)-F'(0)$ likewise have
nonnegative valuation eventually. Therefore $F'(0)\in V$.

Additivity of $\beta$ is immediate. For $r\in R$,
\[
 D(rF'(0))G'(0)-rD(F'(0))G'(0)
   =D(r)F'(0)G'(0)\in V.
\]
Thus $\beta(rF,G)=r\beta(F,G)=\beta(F,rG)$, so it induces
an $R$-linear functional on $A\otimes_RA$ with target $K/V$.
On the displayed tensor its value is
\[
 \bigl(D(cP'(0))P'(0)-D(P'(0))cP'(0)\bigr)+V
   =D(c)P'(0)^2+V,
\]
which is nonzero by assumption. Its polynomial image is zero
by cancellation of the two equal products. This proves the
obstruction without any flatness assumption.
\end{proof}

\begin{saresult}{B.039}{Proved}{An explicit nonfinite-normalization counterexample}
Let $k=\mathbb F_2$, $A_0=k[t]_{(t)}$, $n_0=0$ and
$n_i=(i+1)!$ for $i\ge1$. In $k[[t]]$ set
\[
 z=\sum_{i\ge0}t^{n_i},\qquad
 z_i=\frac{z-\sum_{j<i}t^{n_j}}{t^{n_i}},\qquad
 u=t(z-1),\qquad s_i=(z_i-1)^2.
\]
Define
\[
 C=A_0[u,s_0,s_1,\ldots],\qquad
 M=(t,u)C,\qquad R=C_M.
\]
Then $R$ is a one-dimensional Noetherian local domain with
residue field $k$, fraction field $K=k(t,z)$ and normalization
a discrete valuation ring $V$ with parameter $t$ and residue
field $k$. Its normalization is not a finite $R$-module.
Nevertheless, unlike the finite-conductor one-branch cases,
$\operatorname{Int}(R)$ has a noninjective tensor map.
Explicitly, with $q(X)=X^2-X$ and $c=z-1$,
\[
 cq\in\operatorname{Int}(R),\qquad
 (cq)\otimes q-q\otimes(cq)\ne0,
\]
while its polynomial image is zero.
\end{saresult}
\begin{proof}
This is an explicit characteristic-two specialization of the
Akizuki--Reid construction. The known ring-theoretic conclusions
are stated in Reid, \emph{Akizuki's counterexample}, the two
theorems on page 2 and their proofs on pages 2--4; the opened
author LaTeX source is lines 49--212. We verify the hypotheses
and give the short characteristic-two arguments needed below,
including a direct nonfiniteness proof.

First $z$ is transcendental over $k(t)$. Otherwise let
$0\ne P(t,Z)\in k[t,Z]$ vanish at $z$, with $Z$-degree $d$
and coefficient $t$-degrees at most $h$. Its truncation
$z^{(N)}=\sum_{i\le N}t^{n_i}$ gives
\[
 v_t(P(t,z^{(N)}))\ge n_{N+1},\qquad
 \deg_t P(t,z^{(N)})\le h+d n_N.
\]
For all sufficiently large $N$, $n_{N+1}>h+dn_N$, so the
polynomial $P(t,z^{(N)})$ is zero. The distinct elements
$z^{(N)}\in k(t)$ would then give infinitely many roots of the
nonzero polynomial $P(t,Z)$ over that field, a contradiction.
The gap condition $n_i\ge2n_{i-1}+2$ of Reid's theorem also
holds, including $i=1,2$.

Here are direct checks of the local-ring assertions in this
specialization. Put $B=A_0[z_0,z_1,\ldots]$. The equations
\[
 z_i-1=t^{n_{i+1}-n_i}z_{i+1}
\]
make $tB$ the kernel of the residue map $B\to k$.
The local domain $V=B_{tB}$ has principal maximal ideal $tV$
and $\bigcap_n t^nV=0$, because it embeds in $k[[t]]$.
Every nonzero element therefore equals a power of $t$ times
a unit. Every nonzero ideal has an element of least such power,
so is principal; this proves that $V$ is a discrete valuation
ring, without presupposing Noetherianity of $B$.

In characteristic two,
\[
 s_i=t^{2(n_{i+1}-n_i)}(1+s_{i+1}),\qquad u^2=t^2s_0.
\]
Thus every $s_i$ lies in $tC$, and $C/M=k$.
Iterating the first identity shows that, modulo $t^NC$,
each of finitely many specified $s_i$ is represented by an
element of $A_0$. Using the second identity to reduce powers
of $u$ shows that $C/t^NC$ is generated over $A_0/(t^N)$ by
$1,u$. More precisely, with
$u_r=t(z-\sum_{j\le r}t^{n_j})=t^{n_r+1}(z_r-1)$,
every $f\in C$ can be written
\[
 f=X+Yu_r+t^NZ,\qquad X,Y\in A_0,\quad Z\in C,
\]
for any prescribed $r,N$, since $u-u_r\in A_0$.
The radical of $tC$ is $M$, as $C/tC$ is generated by the
square-zero image of $u$. Hence localization at $M$ does not
change $C/t^NC$, which is a finite $A_0/(t^N)$-module.

For $0\ne f\in M$, let $v_t(f)=a$, choose $N>a$ and then
$r$ with $n_r\ge N$. In the preceding expression $X$ has
valuation $a$. Divide by its unit factor in $A_0$ to write
\[
 f=t^a(1+t^{N-a}Z)+Yu_r.
\]
Squaring in characteristic two and using
$u_r^2=t^{2n_r+2}s_r$ gives
$f^2=t^{2a}w$ with $w\in C\setminus M$.
Thus $fR$ contains $t^{2a}$. The same conclusion holds for
any nonzero element of $R$ after clearing its unit denominator.
Every nonzero ideal consequently contains a power of $t$;
its quotient by that principal ideal is a submodule of a
finite $A_0/(t^N)$-module, and is finitely generated. This proves
that $R$ is Noetherian. Its only nonzero prime is $MR$, so its
dimension is one and its residue field is $k$.

Both $R$ and $V$ have fraction field $K$, since $z=1+u/t$.
Every $z_i$ is integral over $C$, as $z_i^2=1+s_i\in C$.
After first localizing $B$ by $C\setminus M$, all its maximal
ideals lie over $MR$ and contain $t$, so the only one is the
extension of $tB$. This localized ring is therefore $B_{tB}=V$
and is integral over $R$. As $V$ is integrally closed in $K$,
it is the normalization.

The derivation $D=t^{-1}\partial/\partial z$ of $K$ preserves
$R$: it kills $A_0$ and all the squares $s_i$, sends $u$ to $1$,
and extends to the localization by the quotient rule.
If $V$ were a finite $R$-module with generators $v_1,\ldots,v_h$,
then the valuations of $D(V)$ would be uniformly bounded below
by the minimum of $0,v(D(v_1)),\ldots,v(D(v_h))$.
Indeed $D(rv_j)=D(r)v_j+rD(v_j)$ and $D(r),r\in R\subset V$.
But $z_i\in V$ and
\[
 D(z_i)=t^{-n_i-1},
\]
whose valuations are unbounded below. This proves nonfiniteness
directly.

The elements $u_r\in R$ satisfy $D(u_r)=1$ and
$v_t(u_r)=n_{r+1}+1\to\infty$, so B.038 applies.
Also $cMR\subset R$, since $ct=u$ and $cu=t s_0$ both lie in
$R$. As $q(R)\subset MR$ by the residue field $k=\mathbb F_2$,
we have $cq\in\operatorname{Int}(R)$.
Finally $D(c)=1/t\notin V$ and $q'(0)=-1$.
The functional of B.038 takes the displayed tensor to
$1/t+V\ne0$, whereas its polynomial image is zero.
This proves the asserted counterexample with all its hypotheses.
\end{proof}


\begin{saresult}{B.040}{Proved}{Every finite-residue Akizuki--Reid family fails tensor uniqueness}
Let $A_0$ be any discrete valuation ring with uniformizer $t$ and
finite residue field of cardinality $q$. Choose units $a_i\in A_0$
and exponents $n_0=0$, $n_{i+1}\ge2n_i+2$, such that
$z=\sum_i a_it^{n_i}\in\widehat A_0$ is transcendental over
$\operatorname{Frac}(A_0)$. Define $z_i,C,M,R$ by Reid's construction:
\[
 z_i=\frac{z-\sum_{j<i}a_jt^{n_j}}{t^{n_i}},\qquad
 C=A_0[t(z-a_0),(z_i-a_i)^2:i\ge0].
\]
Put
\[
 M=(t,t(z-a_0))C,\qquad R=C_M.
\]
Then $R$ is a one-dimensional Noetherian local domain with
nonfinite discrete valuation normalization $V$, as in Reid's
opened theorem. With $c=z-a_0$ and $P=X^q-X$, one has
$cP\in\operatorname{Int}(R)$ and
\[
 (cP)\otimes P-P\otimes(cP)\ne0,
\]
although its polynomial image is zero. This includes mixed
characteristic; it is not restricted to characteristic two.
\end{saresult}
\begin{proof}
The ring assertions, with exactly the hypotheses just stated,
are Reid's two theorems on page 2, with proofs on pages 2--4.
Set $D=t^{-1}\partial/\partial z$ on
$K=\operatorname{Frac}(A_0)(z)$. This derivation kills $A_0$
and sends $t(z-a_0)$ to $1$. On the remaining generators,
\[
 D((z_i-a_i)^2)=2t^{-n_i-1}(z_i-a_i).
\]
Its valuation is at least
$n_{i+1}-2n_i-1\ge1$; if $2=0$ the value is zero.
Thus $D(C)\subset V$. Denominators from $C\setminus M$ are
units of $V$, so the quotient rule gives $D(R)\subset V$.
The elements
\[
 u_i=t\left(z-\sum_{j\le i}a_jt^{n_j}\right)
      =t^{n_i+1}(z_i-a_i)\in R
\]
satisfy $D(u_i)=1$ and $v_t(u_i)=n_{i+1}+1\to\infty$.
Consequently B.038 applies with its weaker, explicitly stated
hypothesis $D(R)\subset V$.

Writing $u=t(z-a_0)$, one has $ct=u$ and
$cu=t(z-a_0)^2\in C$, so $cMR\subset R$.
The polynomial $P$ has values in $MR$ because the residue field
has $q$ elements. Hence $cP\in\operatorname{Int}(R)$.
Finally $D(c)=1/t$ and $P'(0)=-1$; B.038 takes the tensor to
$1/t+V\ne0$ and proves the assertion.
The same bounded-derivative argument as in B.039 reproves
nonfiniteness of $V/R$, since $D(z_i)=t^{-n_i-1}$ has
unboundedly negative valuation.
\end{proof}


\begin{saresult}{B.041}{Proved}{Every infinite valuation domain and every torsion-free extension}
Let $V$ be an infinite valuation domain with fraction field $K$,
and let $V\subset S$ be any $V$-torsion-free unital extension.
Then every polynomial tensor map for $\Omega(V,S)$ is an
isomorphism, and the unique compatible plethory exists.
More precisely, put $L=S_{\mathrm{rat}}$. If the residue field
of $V$ is infinite or its maximal ideal is nonprincipal, then
$\Omega(V,S)=V[X]$. If its nonzero maximal ideal is principal,
with finite residue field, then
\[
 \Omega(V,S)=
 \begin{cases}
 V[X],&L=V,\\
 \operatorname{Int}(V),&L\ne V,
 \end{cases}
\]
and $\operatorname{Int}(V)$ has an explicit free Newton basis.
This covers valuation domains of arbitrary rank and does not
assume Noetherianity. A finite valuation domain is a finite field
and is governed instead by B.013.
\end{saresult}
\begin{proof}
The field case is B.014, so suppose $V$ is not a field.
B.014 reduces $S$ to the overring $L\subset K$.
An infinite residue field gives $\operatorname{Int}(V)=V[X]$
by the invertible Vandermonde argument; hence it gives the
asserted equality for every coefficient overring.

Suppose the maximal ideal is nonprincipal. The ordered value
group $\Gamma$ has no least positive element: a least positive
value would give a generator of that ideal, and conversely.
For every $\gamma>0$ and every positive integer $N$, there is
$\delta>0$ with $N\delta<\gamma$. To verify this without an
Archimedean assumption, choose $0<\epsilon<\gamma$; any positive
$\epsilon$ has a positive $\eta$ with $2\eta\le\epsilon$.
Indeed choose $0<\xi<\epsilon$ and take $\eta=\xi$ if
$2\xi\le\epsilon$, otherwise take $\eta=\epsilon-\xi$.
Iterate until $2^j\ge N$.

Take $F\in\operatorname{Int}(V)$ of degree $d\ge1$. If one
coefficient had value $-\gamma<0$, choose $a\in V$ with value
$\delta>0$ so small that
\[
 \frac{d(d+1)(d+2)}6\,\delta<\gamma.
\]
The $d+1$ points $0,a,a^2,\ldots,a^d$ are distinct. Their
Vandermonde determinant has value exactly
$d(d+1)(d+2)\delta/6$: the difference of distinct positive
powers of $a$ has the smaller of their two values.
Cramer's rule, with all values $F(a^i)$ in $V$, forces every
coefficient to have value at least the negative of this
quantity, contradicting $-\gamma$. Constant polynomials are
already in $V$. Thus again $\operatorname{Int}(V)=V[X]$.

It remains to suppose the maximal ideal is $\pi V$ and the
residue field has $q$ elements. Choose representatives
$c_0=0,c_1,\ldots,c_{q-1}\in V$. For the base-$q$ expansion
$n=\sum_j n_jq^j$ put $a_n=\sum_jc_{n_j}\pi^j$ and
\[
 w_q(n)=\sum_{j\ge1}\lfloor n/q^j\rfloor,\qquad
 b_n(X)=\frac{\prod_{i<n}(X-a_i)}{\prod_{i<n}(a_n-a_i)}.
\]
The denominator has value exactly $w_q(n)v(\pi)$, by the
first differing digit of each pair. Every $x\in V$ has any
prescribed finite number of successive residue digits, since
its maximal ideal is $\pi V$. Among $a_0,\ldots,a_{n-1}$ at
least $\lfloor n/q^j\rfloor$ share the first $j$ digits of
$x$. Summing these finite lower bounds shows that the
numerator has value at least $w_q(n)v(\pi)$.
This argument uses only finitely many digit levels with
nonzero lower bounds, so it remains valid in a higher-rank
value group. Hence $b_n\in\operatorname{Int}(V)$ and
$b_n(a_n)=1$. Successive evaluation at these nodes gives
\[
 \operatorname{Int}(V)=\bigoplus_{n\ge0}Vb_n.
\]
Every $b_n$ has coefficients in $V[\pi^{-1}]$.

If $L\ne V$, choose $b\in L\setminus V$.
The least positive value $v(\pi)$ satisfies
$-v(b)\ge v(\pi)$. Therefore $(\pi b)^{-1}\in V$ and
$\pi^{-1}=(\pi b)^{-1}b\in L$.
The entire Newton basis belongs to $L[X]$, so
$\Omega(V,L)=\operatorname{Int}(V)$. If $L=V$ the module is
$V[X]$. These modules are free over $V$ in every case;
B.018 gives all tensor isomorphisms and the plethory.
\end{proof}


\begin{saresult}{B.042}{Proved}{Every infinite Krull domain with arbitrary torsion-free coefficients}
Let $R$ be an infinite Krull domain with fraction field $K$, and
let $R\subset S$ be an $R$-torsion-free unital extension.
Then $\Omega(R,S)$ is a projective $R$-module, every
multivariable tensor map is an isomorphism, and the unique
compatible plethory exists. No Noetherian hypothesis is needed.
Here the Krull hypotheses used are precisely
$R=\bigcap_{\operatorname{ht}\mathfrak p=1}R_{\mathfrak p}$,
each $R_{\mathfrak p}$ a discrete valuation ring, and every
nonzero element of $R$ lying in only finitely many such primes.
\end{saresult}
\begin{proof}
The field case is B.014. Put $L=S_{\mathrm{rat}}\subset K$.
We verify the two places where B.020 used Noetherianity and
then complete the fractional-ideal decomposition directly.

Every height-one maximal ideal $\mathfrak p$ is invertible.
Choose $b\in R$ with $v_{\mathfrak p}(b)=1$. Besides
$\mathfrak p$, only finitely many height-one primes
$\mathfrak q_i$ contain $b$. Choose
$a_i\in\mathfrak q_i^{v_{\mathfrak q_i}(b)}\setminus\mathfrak p$
and put $a=\prod_i a_i$. Such choices exist because distinct
height-one primes are incomparable; powers of an element of
$\mathfrak q_i\setminus\mathfrak p$ suffice.
Then $x=a/b$ has valuation $-1$ at $\mathfrak p$ and
nonnegative valuation at every other height-one prime.
Consequently $x\mathfrak p\subset R$, by the Krull intersection
property. The ideal $\mathfrak p\mathfrak p^{-1}$ contains
$\mathfrak p$ and $xb=a\notin\mathfrak p$, so equals $R$.
This proves invertibility rather than assuming finite generation.

For each fixed prime power $q$, choose $r_q\in R$ with
$r_q^q-r_q\ne0$; this is possible since $R$ is infinite.
Every height-one maximal ideal with residue field of size $q$
contains this nonzero element, so there are only finitely many
such ideals. In particular the following product is finite:
\[
 I_n=\prod_{\substack{\operatorname{ht}\mathfrak p=1,\,
                       \mathfrak p\text{ maximal}\\
                       |R/\mathfrak p|=q<\infty}}
           \mathfrak p^{-w_q(n)},\qquad
 w_q(n)=\sum_{j\ge1}\lfloor n/q^j\rfloor.
\]
All its factors, and hence $I_n$, are invertible fractional ideals.
An invertible fractional ideal is the intersection of its
height-one localizations: if $x$ belongs to those localizations,
then $xI_n^{-1}$ is integral at every height-one prime and hence
lies in $R$, which implies $x\in I_n$.

The digit Newton argument at a finite-residue height-one maximal
ideal shows that the leading coefficient of an integer-valued
polynomial of degree at most $n$ has valuation at least
$-w_q(n)$. The necessary nodes can be chosen in $R$ by successive
residue digits and a valuation-one element in $R$.
At every other height-one prime, the image of $R$ in its residue
field is infinite, so Vandermonde interpolation gives
nonnegative valuations of all coefficients. The leading
coefficient ideal is therefore contained in $I_n$.

For each prime occurring in the finite product, choose its monic
local digit polynomial $Q_{n,\mathfrak p}\in R[X]$.
The Chinese remainder theorem for the pairwise comaximal
powers $\mathfrak p^{w_q(n)}$ gives a monic polynomial
$P_n\in R[X]$ of degree $n$, congruent coefficientwise to
$Q_{n,\mathfrak p}$ modulo that power at each such prime.
For $c\in I_n$ and $r\in R$, the product $cP_n(r)$ is integral
at every height-one prime, and so belongs to $R$.
Subtracting its leading term successively proves
\[
 \operatorname{Int}(R)=\bigoplus_{n\ge0}I_nP_n,\qquad
 \Omega(R,L)=\bigoplus_{n\ge0}(I_n\cap L)P_n.
\]
The second equality uses that the $P_n$ are monic with
coefficients in $R$, so descending elimination preserves $L[X]$.

Put $J_n=I_n\cap L$ and let $E_n$ be the finite set of primes
occurring in $I_n$. Localization commutes with this finite
intersection of submodules of $K$. At a maximal ideal outside
$E_n$, $J_n$ localizes to $R_{\mathfrak m}$. At
$\mathfrak p\in E_n$, the overring $L_{\mathfrak p}$ of the
rank-one discrete valuation ring $R_{\mathfrak p}$ is either
$R_{\mathfrak p}$ or $K$. Thus localization at every maximal
ideal proves the exact formula
\[
 J_n=\prod_{\substack{\mathfrak p\in E_n\\L_{\mathfrak p}=K}}
                 \mathfrak p^{-w_{|R/\mathfrak p|}(n)}.
\]
No finite-generation assertion about an arbitrary intersection
has been used: equality is proved by localization against the
displayed finite invertible product. Each $J_n$ is therefore
projective, and their direct sum is projective. B.019 proves
all tensor and plethory assertions.
\end{proof}

\begin{saresult}{B.043}{Proved}{Infinite-residue conductor cusps and all their coefficients}
Let $k=\mathbb F_q$ be finite, let $l\supset k$ be any infinite
field, and put
\[
 V=l[[t]],\qquad R=k+t^aV\quad(a\ge1),\qquad K=l((t)).
\]
For every $R$-torsion-free unital extension $R\subset S$, put
$L=S_{\mathrm{rat}}$. Then the polynomial coaddition and
comultiplication tensors are uniquely determined if and only
if $L=R$. In particular fraction-field coefficients fail.
For $A=\operatorname{Int}(R)$, $s\in l\setminus k$,
$Q(X)=X^q-X$ and $F(X)=Q(X)^2/t^a$, an explicit kernel is
\[
 (sF)\otimes F-F\otimes(sF)\ne0.
\]
The ring $R$ is a non-Noetherian local domain with maximal ideal
$t^aV$, residue field $k$, and fraction field $K$.
\end{saresult}
\begin{proof}
The ring assertions except non-Noetherianity follow from its
pullback description and from $t^aV\subset R$: an element with
nonzero constant in $k$ has inverse in $k+t^aV$, and division
by a power of $t^a$ gives every element of $K$.
Its maximal ideal $\mathfrak m=t^aV$ has
$\mathfrak m/\mathfrak m^2$ of infinite dimension over $k$,
since it contains $t^al$ and $[l:k]=\infty$.
Thus $\mathfrak m$ is not finitely generated.

We first treat coefficient field $K$ and construct a functional
which proves that the proposed tensor is nonzero.
For $G\in A$, the polynomial $G(t^aY)$ takes $V$ into
$R\subset V$. Since $V$ has infinite residue field,
$\operatorname{Int}(V)=V[Y]$ by Vandermonde interpolation.
Its reduction modulo $t$ is a polynomial over $l$ taking all
of $l$ into $k$. Such a polynomial is constant: the product of
its differences from all elements of $k$ vanishes at every
point of the infinite field $l$, hence is the zero polynomial;
one factor must be zero. Its constant is $G(0)\bmod t$.
It follows that
\[
 G(t^ay)-G(0)\in R\cap tV=t^aV\quad(y\in V).
\]
Apply the same Vandermonde argument to its quotient by $t^a$.
This defines an $R$-linear map, with $R$ acting through $k$,
\[
 h:A\longrightarrow l[Y],\qquad
 h(G)=\left(\frac{G(t^aY)-G(0)}{t^a}\right)\bmod t.
\]
In particular $h$ kills $\mathfrak m A$.

For $r\in R$ the value $Q(r)$ lies in $t^aV$, so both $F$ and
$sF$ take $R$ into $t^aV\subset R$. Direct expansion gives
\[
 h(F)=Y^2,\qquad h(sF)=sY^2.
\]
The image of the proposed tensor under $h\otimes h$ is
\[
 sY^2\otimes Z^2-Y^2\otimes sZ^2
       \in l[Y]\otimes_k l[Z],
\]
which is nonzero. For example choose a $k$-linear functional
$\ell:l\to k$ with $\ell(1)=0$, $\ell(s)=1$, and apply it
to the first $Y^2$ coefficient and take the second $Z^2$
coefficient. The result is $1$. The polynomial image of the
original tensor is zero by cancellation of its two products.

For completeness consider all coefficient overrings. Any
$L\not\subset V$ is $K$: if $v(b)<0$, choose $m$ with $v(b^{-m})\ge a$; then
$c=b^{-m}\in t^aV\subset R$ and $c^{-1}\in L$. Every element
of $K$ is sent into $R$ by a high power of $c$.
Suppose now $R\subsetneq L\subset V$. The conductor contains
$\mathfrak m$, and $Q(R)\subset\mathfrak m$.
Monic division gives a split summand $LQ$ in $\Omega(R,L)$.
If the second tensor map were injective, multiplication
$L\otimes_RL\to L$ would be an isomorphism. Modulo
$\mathfrak m$, it would give injective multiplication on the
nonzero $k$-algebra $C=L/\mathfrak mL=L/\mathfrak m$.
Choose a $k$-linear functional taking $1$ to $1$.
The relation $c\otimes1=1\otimes c$, forced by injectivity,
then implies that every $c\in C$ belongs to $k$.
Hence $C=k$ and $L=k+\mathfrak m=R$, a contradiction.
This argument does not use a finite-module Nakayama lemma:
$L$ can be infinite over $R$.
The only positive case is therefore $L=R$, which gives the
ordinary polynomial plethory. B.014 transfers the result to
arbitrary torsion-free $S$.
\end{proof}


\begin{saresult}{B.044}{Proved}{Integer-valued polynomials commute with one-dimensional finite-residue completion}
Let $(R,\mathfrak m)$ be a one-dimensional Noetherian local
domain with finite residue field, let $K=\operatorname{Frac}(R)$,
and put $A=\widehat R$ and $Q=K\otimes_R A$.
For every number of variables $n$ there is a canonical isomorphism
\[
 \operatorname{Int}(R^n,R)\otimes_R A
       \simeq \operatorname{Int}(A^n,A)\subset Q[X_1,\ldots,X_n].
\]
It identifies the tensor maps after faithfully flat base change.
In particular a noninjective second tensor map for $A\subset Q$
implies a noninjective second tensor map for $R\subset K$.
\end{saresult}
\begin{proof}
Fix bounds on the degree in each variable and let $N$ be the
number of allowed monomials. Choose a rectangular grid of
distinct points in $R$. The product of its Vandermonde
determinants gives a nonzero $d\in R$ such that every
integer-valued polynomial within these bounds has all
coefficients in $d^{-1}R$. The same grid over $A$ shows that
all coefficients there lie in $d^{-1}A$: flatness makes $d$
a nonzerodivisor of $A$.

The ring $R/dR$ is a finite ring, because it is Artinian local
with finite residue field. Choose a finite set $E\subset R$
of representatives. For $f\in d^{-1}R[X_1,\ldots,X_n]$,
its value modulo $R$ depends only on each input modulo $dR$.
Indeed $df$ has coefficients in $R$, and changing a coordinate
by a multiple of $d$ changes its numerator by a multiple of
$d$. Consequently the integer-valued module in these bounds
is the kernel of the finite evaluation map
\[
 (d^{-1}R)^N\longrightarrow (d^{-1}R/R)^{E^n}.
\]
The natural map $R/dR\to A/dA$ is an isomorphism: the source
has finite length and hence is already complete. Thus exactly
the same finite tests describe the bounded-degree module over
$A$. Flat base change commutes with this kernel and identifies
the two modules. Taking the filtered union of all degree
bounds proves the assertion. Polynomial multiplication in
disjoint variables commutes with these identifications, so
the tensor maps correspond as stated. Faithful flatness
detects nonzero kernels.

For clarity, $Q$ is the total quotient ring of $A$. Choose
$0\ne t\in\mathfrak m$. Every $0\ne r\in R$ has
$t^M\in rR$ for some $M$, so $Q=A[1/t]$. This is a
zero-dimensional Noetherian ring: $A/tA$ is Artinian, and
completion has dimension one. It is Artinian and is a
finite product of local Artinian rings. Every nonzerodivisor
of $A$ becomes a unit in this product, so $A[1/t]$ is indeed
the total quotient ring.
The completion facts invoked here are the flatness and exactness
of finite-module completion (Stacks, Lemma 10.97.2, tag 00MB)
and its local faithful-flat consequence.
\end{proof}

\begin{saresult}{B.045}{Proved}{A generic nilpotent forces a tensor kernel after completion}
Let $(A,\mathfrak m)$ be a one-dimensional complete
Noetherian local ring with finite residue field. Assume that
$A$ is Cohen--Macaulay and that either $A$ has characteristic
$p>0$, or $A$ has characteristic zero and $p$ is a
nonzerodivisor, where $p$ is its residue characteristic.
If $A$ is not reduced, then
\[
 \operatorname{Int}(A)\otimes_A\operatorname{Int}(A)
       \longrightarrow\operatorname{Int}(A^2)
\]
is not injective. Here coefficients belong to the total
quotient ring of $A$. The proof constructs a nonzero tensor
and an explicit balanced linear functional detecting it.
\end{saresult}
\begin{proof}
We first construct the finite lattice in which the argument
will take place, including the needed compatible coefficient
field. There is a complete discrete valuation subring
$V_0\subset A$ such that $A$ is finite and torsion-free over
$V_0$, and such that some nonreduced local factor $Q_0$ of
$Q(A)$ has residue field $F$ separable over
$K_0=\operatorname{Frac}(V_0)$.

In mixed characteristic take the Cohen coefficient ring.
It is a complete discrete valuation ring with uniformizer
$p$, and $p$ is a parameter. The finite-module assertion
follows by lifting finitely many generators modulo $p$ and
successively approximating in the $p$-adic topology;
$A$ is $p$-adically complete because $pA$ is
$\mathfrak m$-primary. Separability holds in characteristic
zero.

Here is the extra parameter choice needed in characteristic
$p$. The roots of $X^{|A/\mathfrak m|}-X$ give a coefficient
field $k\subset A$ (or use the Cohen structure theorem).
Choose a minimal prime $\mathfrak p$ corresponding to a
nonreduced generic factor. Such a factor exists because a
parameter is a nonzerodivisor, so a nonzero nilpotent remains
nonzero in $Q(A)$. The complete one-dimensional domain
$C=A/\mathfrak p$ has finite normalization, a complete
discrete valuation ring $V=l[[u]]$ with $l/k$ finite.
Here finiteness follows from the Nagata property of complete
Noetherian local rings; a finite domain over a complete local
ring is local, because a finite complete algebra decomposes
as a product of complete local algebras. A nonzero conductor
then gives $u^N\in C$ for all sufficiently large $N$.
Choose such an $N$ with $p\nmid N$ and lift $u^N$ to
$t_0\in\mathfrak m$.

We can adjust this lift to a parameter without changing its
image in $C$. Choose
$a\in\mathfrak p$ outside every other minimal prime, by
prime avoidance. At each other minimal prime, the nonzero
images of $a,a^2,a^3,\ldots$ are distinct: $a$ lies in
the maximal ideal of the corresponding local domain.
Thus $t_0+a^j$ belongs to each such prime for at most one
$j$. Choosing $j$ outside this finite exceptional set gives
a parameter $t$ mapping to $u^N$. If there is no other
minimal prime, use $t=t_0$.
Set $V_0=k[[t]]\subset A$. This map is injective, since
a nonzero power series is $t^r$ times a unit and $t$ is a
nonzerodivisor. Successive approximation modulo $t^m$ shows
that $A$ is finite over $V_0$. In the chosen factor the
residue field is $F=l((u))$, a finite separable extension
of $k((u^N))$ since $p\nmid N$.

Since $F/K_0$ is finite separable, it embeds in $Q_0$ as
a coefficient field compatibly with $K_0$. To see this
directly, choose a primitive element of $F/K_0$ and lift
it to $Q_0$. Its separable minimal polynomial has invertible
derivative. Newton correction modulo successive powers of
the nilradical, which is nilpotent, gives an actual root
lifting the primitive element. Its $K_0$-algebra is the
required copy of $F$. Write
\[
 \rho:Q_0\longrightarrow F,\qquad N=\ker\rho,
 \qquad N^h=0\quad(h\ge2).
\]
Let $V\subset F$ be the integral closure of $V_0$ in $F$.
It is a finite complete discrete valuation ring, with
uniformizer $u$ and finite residue field. The same notation
$V=l[[u]]$ applies in the equal-characteristic construction.
The ring
\[
 O=V[\operatorname{pr}_0(A)]\subset Q_0
\]
is finite and torsion-free over $V$, contains $V$, satisfies
$\rho(O)=V$, and spans $Q_0$ over $F$.
Finiteness follows because $\operatorname{pr}_0(A)$ is
finite over $V_0$ and hence integral over $V$.
Let $e$ denote the idempotent of the selected factor in
$Q(A)$. Since $eO$ is a finite $V_0$-module in $Q(A)$,
clearing finitely many denominators gives an integer $c\ge0$
such that
\[
                         u^c eO\subset A.                 \tag{B45.1}
\]
Indeed a sufficiently high power of the uniformizer of
$V_0$ works, and in $V$ that power is a unit times a
power of $u$; the unit preserves $O$.

Put $J=N\cap O$. There is a primitive element
$z_0\in J\cap\operatorname{Ann}_{Q_0}(N)$: start with a
nonzero element in the last nonzero power of $N$, multiply
it into $O$, and divide by the greatest possible power of
$u$ in the finite free $V$-module $O$. Choose a
$V$-linear map $\ell:O\to V$ with $\ell(z_0)=1$ and
extend it $F$-linearly to $Q_0$. Such a map exists because
a primitive vector in a finite free module over a discrete
valuation ring extends to a basis.

There is an integer $b\ge0$, independent of
$G\in\operatorname{Int}(A)$, such that, with
$x_G=\operatorname{pr}_0(G'(0))$,
\[
                  Jx_G\subset u^{-b}O,\qquad
                  \rho(x_G)\in u^{-b}V.                   \tag{B45.2}
\]
Here is a finite interpolation proof. Choose finitely many
$V$-generators $n_1,\ldots,n_s$ of $J$. By (B45.1),
$u^c e n_j\in A$. For any $h$ fixed distinct scalars
$s_1,\ldots,s_h\in V_0$, the values
$G(s_i u^c e n_j)$ lie in $A$ and their selected components
lie in $O$. As a polynomial in the scalar $s$, this
evaluation has degree at most $h-1$, because $n_j^h=0$.
Inverting this one fixed Vandermonde matrix bounds its
linear coefficient $u^c n_jx_G$ in a fixed fractional
lattice $u^{-d}O$. Increasing $b$ gives the first inclusion
for every generator, hence all of $J$. Since $z_0N=0$,
$z_0x_G=z_0\rho(x_G)$. Apply $\ell$ to obtain the second
inclusion from the first.

Write $\lambda(G)=\rho(x_G)$. Define
\[
 \beta(G,H)=\ell(x_G)\lambda(H)\pmod {u^{-2b}V}
                         \quad\hbox{in }F/u^{-2b}V.
                                                               \tag{B45.3}
\]
This additive pairing is $A$-balanced. If $a\in A$ and
$a_0=\operatorname{pr}_0(a)$, the difference between
$\beta(aG,H)$ and $\beta(G,aH)$ equals
\[
 \ell\bigl((a_0-\rho(a_0))x_G\bigr)\lambda(H).
\]
The first factor lies in $u^{-b}V$ by (B45.2), since
$a_0-\rho(a_0)\in J$, and the second does also.
Their product therefore vanishes in the indicated quotient.
The quotient is an $A$-module through $\rho\operatorname{pr}_0$,
whose image is in $V$.

Choose an integer $r>2c+2b$. We construct a polynomial
$g_r\in F[X]$ with
\[
 g_r(0)=0,\quad g_r'(0)=(-u^{-1})^r,\quad
 g_r(V)\subset V,\quad g_r(O)\subset V+u^{-r}J.             \tag{B45.4}
\]
In characteristic $p$, choose a power $D$ of the residue
cardinality with $D\ge h$ and set
$\delta(X)=(X^D-X)/u$, $g_r=\delta^{\circ r}$.
The Frobenius identity and $N^h=0$ give
$\delta(a+n)=\delta(a)-n/u$ for $a\in V,n\in N$;
all assertions follow.

In mixed characteristic choose a power $D$ of the residue
cardinality so large that $D\ge h$ and
\[
 v(D)\ge (r-1)(h-2)+\max_{1\le j<h}v(j).
\]
Again use $g_r=\delta^{\circ r}$. The elementary identity
$\binom Dj=(D/j)\binom{D-1}{j-1}$ gives
$v(\binom Dj)\ge v(D)-v(j)$.
For $0\le i<r$, $a\in V$ and $n\in u^{-i}J$, expand
\[
 \delta(a+n)=\delta(a)+\frac{Da^{D-1}-1}{u}n
          +\sum_{2\le j<h}\frac{\binom Dj a^{D-j}}u n^j.
\]
The linear term belongs to $u^{-(i+1)}J$, and so does
each higher term: its valuation bound is
$v(\binom Dj)-ij-1\ge-i-1$.
Thus $\delta(V+u^{-i}J)\subset V+u^{-(i+1)}J$.
Induction proves (B45.4); its derivative follows from
$\delta(0)=0$ and $\delta'(0)=-u^{-1}$.

Extend polynomials from $Q_0$ to $Q(A)$ by zero on the
other factors, and put
\[
       F_r=u^c e z_0g_r,\qquad H_r=u^{r+c}eg_r.
\]
Both are in $\operatorname{Int}(A)$. For $H_r$ this is
(B45.1) and (B45.4). For $F_r$, multiplication by $z_0$
kills the nilpotent part of $g_r(a)$, leaving
$u^cz_0g_r(\rho(a))\in u^cO$.
The polynomial image of
\[
                   F_r\otimes H_r-H_r\otimes F_r
\]
is zero, because $F_r=u^{-r}z_0H_r$ in $Q(A)[X]$.
Its value under (B45.3) is $u^{2c-r}$: the second term
vanishes since $\rho(z_0)=0$, and the two signs cancel
in the first. This is nonzero modulo $u^{-2b}V$ by the
choice of $r$. The required tensor kernel is proved.

The structural inputs were opened in Stacks, Theorem 10.160.8
(Cohen structure, tag 032A), Lemma 10.160.2 (complete finite
algebras), and Lemma 10.162.8 (complete local rings are Nagata,
tag 032W). The lattice and obstruction arguments above are
given in full and do not invoke flatness of integer-valued
polynomial rings.
\end{proof}

\begin{saresult}{B.046}{Proved}{Complete one-dimensional fraction-field classification without finite normalization}
Let $R$ be an infinite one-dimensional Noetherian domain,
and put $K=\operatorname{Frac}(R)$.
The compatible coaddition and comultiplication tensors for
$\operatorname{Int}(R)$ are uniquely determined if and only
if, at each maximal ideal with finite residue field,
the completion $\widehat{R_{\mathfrak m}}$ is a domain.
Equivalently, at each such ideal the normalization of
$R_{\mathfrak m}$ is finite and local. In every positive
case all multivariable tensor maps are isomorphisms and
the compatible plethory is unique.
\end{saresult}
\begin{proof}
By B.033 the tensor question localizes, so fix a
finite-residue maximal ideal and write its local ring
as $R$. If its completion is nonreduced, B.045 applies:
the completion is one-dimensional Cohen--Macaulay, and
in mixed characteristic the nonzero element $p\in R$
remains a nonzerodivisor by flatness. B.044 then proves
failure of tensor injectivity over $R$.

If the completion is reduced, $R$ is analytically
unramified and its normalization $T$ is finite
(Stacks, Lemma 10.162.10(4), tag 032Y, whose full proof
was opened). B.025 proves that the tensors work exactly
when $T$ is local. This is equivalent to the completion
of $R$ being a domain. In one direction, when $T$ is
local it is a discrete valuation ring, and exact
completion embeds $\widehat R$ in the complete discrete
valuation ring $\widehat T$. Conversely, if
$\widehat R$ is a domain, flatness embeds
$T\otimes_R\widehat R$ into
$K\otimes_R\widehat R$, which is a field. But the
completion of the finite semilocal algebra $T$ is the
product of its nonzero completions at its maximal
ideals. A subring of a field cannot have the nontrivial
idempotents of a product with more than one factor.
Thus $T$ is local.

At maximal ideals of infinite residue field the ordinary
Vandermonde argument gives the polynomial ring, as
already proved. In the positive finite-residue case,
B.024 gives every multivariable tensor isomorphism.
Testing kernels and cokernels at every maximal ideal
proves the global assertion, and B.009 supplies the
unique plethory compatibility.
\end{proof}
